% Additional tables relocated from the main manuscript for length.
\section{Additional main-text tables}
\label{app:main_tables}

Table~\ref{tab:descriptives} reports answer-level descriptive statistics.
Table~\ref{tab:baseline} reports baseline Post$\times$ASI associations.
Table~\ref{tab:ftests} reports joint event-study $F$-tests.
Table~\ref{tab:hsummary} summarizes primary and secondary tests.

\begin{table*}[t]
\centering
\caption{Descriptive statistics (answer-level sample, $N = 4{,}215{,}082$)}
\label{tab:descriptives}
\scriptsize
\setlength{\tabcolsep}{2pt}
\begin{tabular}{lcccccccc}
\toprule
Variable & $N$ & Mean & SD & P25 & Median & P75 & Min & Max \\
\midrule
Net vote score & 4,215,082 & 3.14 & 8.42 & 0 & 1 & 4 & $-12$ & 1,405 \\
Accepted answer (binary) & 4,215,082 & 0.42 & 0.49 & 0 & 0 & 1 & 0 & 1 \\
Answer comments & 4,215,082 & 1.85 & 3.21 & 0 & 1 & 2 & 0 & 84 \\
Informational index ($\mathrm{ICI}_a$) & 4,215,082 & 0.64 & 0.22 & 0.48 & 0.65 & 0.82 & 0 & 1 \\
Epistemic index ($\mathrm{ECI}_a$) & 4,215,082 & 0.38 & 0.25 & 0.18 & 0.34 & 0.55 & 0 & 1 \\
AI answerability ($\mathrm{ASI}_q$) & 1,842,110 & 0.55 & 0.28 & 0.32 & 0.56 & 0.78 & 0 & 1 \\
Complexity & 1,842,110 & 0.41 & 0.30 & 0.15 & 0.38 & 0.62 & 0 & 1 \\
Pre-event reputation (raw source variable) & 354,209 & 4,520 & 12,450 & 15 & 410 & 3,850 & 1 & 1.2M \\
Tag volatility & 14,502 & 0.22 & 0.18 & 0.08 & 0.19 & 0.31 & 0 & 0.88 \\
\bottomrule
\end{tabular}
\end{table*}

\begin{table}[t]
\centering
\caption{Joint $F$-tests for event-study dynamics and pre-trends}
\label{tab:ftests}
\scriptsize
\setlength{\tabcolsep}{4pt}
\begin{tabular}{llccc}
\toprule
Test & Coefficients tested & Num.\ df & $F$ & $p$ \\
\midrule
Local pre-trend leads & $k=-6$ to $-2$ & 5 & 0.86 & 0.51 \\
Code-LLM leads & $k^{C}=-6$ to $-1$ & 6 & 0.71 & 0.64 \\
Code-LLM lags & $k^{C}=0$ to $+5$ & 6 & 1.27 & 0.27 \\
Eight pre-window bins & Table~\ref{tab:prewindow} & 8 & 1.14 & 0.33 \\
Post-arrival months & $k=0$ to $+12$ & 13 & 4.38 & $<$0.001 \\
\bottomrule
\multicolumn{5}{@{}p{\textwidth}@{}}{\scriptsize Joint $F$-tests evaluate
the indicated event-study bins. Each row names the coefficients in the
contrast. The plotted window and the tested window are not the same.
$p$-values are computed using cluster-robust standard errors.}
\end{tabular}
\end{table}

\begin{table*}[t]
\centering
\caption{Primary tests and bounded extensions}
\label{tab:hsummary}
\scriptsize
\setlength{\tabcolsep}{3pt}
\begin{tabular}{p{0.06\textwidth}p{0.25\textwidth}p{0.18\textwidth}p{0.10\textwidth}p{0.28\textwidth}}
\toprule
Test & Prediction & Estimate (95\% CI) & Raw $p$ & Evidentiary status \\
\midrule
H1 & Informational relative value declines & $-0.153$ $[-0.224,-0.082]$ & 0.00002 & Core; Bonferroni(12) $p=0.00024$ \\
H2 & Epistemic relative value rises & $0.228$ $[0.148,0.308]$ & $<$0.0001 & Core; remains below $0.05$ under Bonferroni(12) \\
H3 & Informational association more negative at high ASI & $-0.139$ $[-0.227,-0.051]$ & 0.0020 & Core; Bonferroni(12) $p=0.024$ \\
H4 & Epistemic association more positive at high ASI & $0.176$ $[0.074,0.278]$ & 0.00071 & Core; Bonferroni(12) $p=0.0086$ \\
H5 & Higher-status users shift toward epistemic work & $0.058$ $[0.034,0.082]$ & $<$0.0001 & Secondary; continuing contributors \\
H6 & Epistemic premium stronger under volatility & $0.079$ $[0.006,0.152]$ & 0.034 & Exploratory; Bonferroni(8) $p=0.272$ \\
H7 & Higher-status users shift toward lower-ASI questions & $-0.082$ $[-0.119,-0.045]$ & $<$0.0001 & Secondary; continuing contributors \\
\bottomrule
\multicolumn{5}{@{}p{\textwidth}@{}}{\scriptsize H1--H4 raw $p$-values are two-sided normal probabilities for the printed coefficient divided by its standard error. H1 is reported as $0.00002$ from $t=-4.25$.}
\end{tabular}
\end{table*}

\begin{table*}[t]
\centering
\caption{Baseline differential post-period associations by AI answerability}
\label{tab:baseline}
\scriptsize
\begin{tabular}{lcccc}
\toprule
& (1) Net vote score (OLS) & (2) Acceptance & (3) Future edits & (4) Survival \\
\midrule
$\mathrm{Post}_t$ & Absorbed & Absorbed & Absorbed & Absorbed \\
$\mathrm{ASI}_q$ & $0.132^{***}$ & $0.032^{*}$ & $0.088^{***}$ & $0.105^{***}$ \\
& (0.033) & (0.014) & (0.021) & (0.028) \\
$\mathrm{Post}_t\times \mathrm{ASI}_q$ & $-0.251^{***}$ & $-0.052^{**}$ & $-0.135^{***}$ & $-0.188^{***}$ \\
& (0.039) & (0.018) & (0.033) & (0.042) \\
Controls/FEs & Yes & Yes & Yes & Yes \\
Observations & 4,215,082 & 4,215,082 & 4,215,082 & 4,215,082 \\
$R^2$ / Pseudo $R^2$ & 0.109 & 0.084 & 0.095 & 0.108 \\
\bottomrule
\multicolumn{5}{@{}p{\textwidth}@{}}{\scriptsize Note: In answer-level
rows, $q=q(a)$. Standard errors are clustered at the tag level.
Post is absorbed by the saturated month fixed effects and is not
separately identified in columns that include month fixed effects.
The interaction terms remain identified as post-period slope
differences relative to the pre-period.
$^{*}p<0.05$, $^{**}p<0.01$, $^{***}p<0.001$.}
\end{tabular}
\end{table*}

% Online appendix / supplementary tables for IJIMDI submission.
% Extracted verbatim from the original manuscript; empirical values unchanged.
% Included from the main manuscript via % Additional tables relocated from the main manuscript for length.
\section{Additional main-text tables}
\label{app:main_tables}

Table~\ref{tab:descriptives} reports answer-level descriptive statistics.
Table~\ref{tab:baseline} reports baseline Post$\times$ASI associations.
Table~\ref{tab:ftests} reports joint event-study $F$-tests.
Table~\ref{tab:hsummary} summarizes primary and secondary tests.

\begin{table*}[t]
\centering
\caption{Descriptive statistics (answer-level sample, $N = 4{,}215{,}082$)}
\label{tab:descriptives}
\scriptsize
\setlength{\tabcolsep}{2pt}
\begin{tabular}{lcccccccc}
\toprule
Variable & $N$ & Mean & SD & P25 & Median & P75 & Min & Max \\
\midrule
Net vote score & 4,215,082 & 3.14 & 8.42 & 0 & 1 & 4 & $-12$ & 1,405 \\
Accepted answer (binary) & 4,215,082 & 0.42 & 0.49 & 0 & 0 & 1 & 0 & 1 \\
Answer comments & 4,215,082 & 1.85 & 3.21 & 0 & 1 & 2 & 0 & 84 \\
Informational index ($\mathrm{ICI}_a$) & 4,215,082 & 0.64 & 0.22 & 0.48 & 0.65 & 0.82 & 0 & 1 \\
Epistemic index ($\mathrm{ECI}_a$) & 4,215,082 & 0.38 & 0.25 & 0.18 & 0.34 & 0.55 & 0 & 1 \\
AI answerability ($\mathrm{ASI}_q$) & 1,842,110 & 0.55 & 0.28 & 0.32 & 0.56 & 0.78 & 0 & 1 \\
Complexity & 1,842,110 & 0.41 & 0.30 & 0.15 & 0.38 & 0.62 & 0 & 1 \\
Pre-event reputation (raw source variable) & 354,209 & 4,520 & 12,450 & 15 & 410 & 3,850 & 1 & 1.2M \\
Tag volatility & 14,502 & 0.22 & 0.18 & 0.08 & 0.19 & 0.31 & 0 & 0.88 \\
\bottomrule
\end{tabular}
\end{table*}

\begin{table}[t]
\centering
\caption{Joint $F$-tests for event-study dynamics and pre-trends}
\label{tab:ftests}
\scriptsize
\setlength{\tabcolsep}{4pt}
\begin{tabular}{llccc}
\toprule
Test & Coefficients tested & Num.\ df & $F$ & $p$ \\
\midrule
Local pre-trend leads & $k=-6$ to $-2$ & 5 & 0.86 & 0.51 \\
Code-LLM leads & $k^{C}=-6$ to $-1$ & 6 & 0.71 & 0.64 \\
Code-LLM lags & $k^{C}=0$ to $+5$ & 6 & 1.27 & 0.27 \\
Eight pre-window bins & Table~\ref{tab:prewindow} & 8 & 1.14 & 0.33 \\
Post-arrival months & $k=0$ to $+12$ & 13 & 4.38 & $<$0.001 \\
\bottomrule
\multicolumn{5}{@{}p{\textwidth}@{}}{\scriptsize Joint $F$-tests evaluate
the indicated event-study bins. Each row names the coefficients in the
contrast. The plotted window and the tested window are not the same.
$p$-values are computed using cluster-robust standard errors.}
\end{tabular}
\end{table}

\begin{table*}[t]
\centering
\caption{Primary tests and bounded extensions}
\label{tab:hsummary}
\scriptsize
\setlength{\tabcolsep}{3pt}
\begin{tabular}{p{0.06\textwidth}p{0.25\textwidth}p{0.18\textwidth}p{0.10\textwidth}p{0.28\textwidth}}
\toprule
Test & Prediction & Estimate (95\% CI) & Raw $p$ & Evidentiary status \\
\midrule
H1 & Informational relative value declines & $-0.153$ $[-0.224,-0.082]$ & 0.00002 & Core; Bonferroni(12) $p=0.00024$ \\
H2 & Epistemic relative value rises & $0.228$ $[0.148,0.308]$ & $<$0.0001 & Core; remains below $0.05$ under Bonferroni(12) \\
H3 & Informational association more negative at high ASI & $-0.139$ $[-0.227,-0.051]$ & 0.0020 & Core; Bonferroni(12) $p=0.024$ \\
H4 & Epistemic association more positive at high ASI & $0.176$ $[0.074,0.278]$ & 0.00071 & Core; Bonferroni(12) $p=0.0086$ \\
H5 & Higher-status users shift toward epistemic work & $0.058$ $[0.034,0.082]$ & $<$0.0001 & Secondary; continuing contributors \\
H6 & Epistemic premium stronger under volatility & $0.079$ $[0.006,0.152]$ & 0.034 & Exploratory; Bonferroni(8) $p=0.272$ \\
H7 & Higher-status users shift toward lower-ASI questions & $-0.082$ $[-0.119,-0.045]$ & $<$0.0001 & Secondary; continuing contributors \\
\bottomrule
\multicolumn{5}{@{}p{\textwidth}@{}}{\scriptsize H1--H4 raw $p$-values are two-sided normal probabilities for the printed coefficient divided by its standard error. H1 is reported as $0.00002$ from $t=-4.25$.}
\end{tabular}
\end{table*}

\begin{table*}[t]
\centering
\caption{Baseline differential post-period associations by AI answerability}
\label{tab:baseline}
\scriptsize
\begin{tabular}{lcccc}
\toprule
& (1) Net vote score (OLS) & (2) Acceptance & (3) Future edits & (4) Survival \\
\midrule
$\mathrm{Post}_t$ & Absorbed & Absorbed & Absorbed & Absorbed \\
$\mathrm{ASI}_q$ & $0.132^{***}$ & $0.032^{*}$ & $0.088^{***}$ & $0.105^{***}$ \\
& (0.033) & (0.014) & (0.021) & (0.028) \\
$\mathrm{Post}_t\times \mathrm{ASI}_q$ & $-0.251^{***}$ & $-0.052^{**}$ & $-0.135^{***}$ & $-0.188^{***}$ \\
& (0.039) & (0.018) & (0.033) & (0.042) \\
Controls/FEs & Yes & Yes & Yes & Yes \\
Observations & 4,215,082 & 4,215,082 & 4,215,082 & 4,215,082 \\
$R^2$ / Pseudo $R^2$ & 0.109 & 0.084 & 0.095 & 0.108 \\
\bottomrule
\multicolumn{5}{@{}p{\textwidth}@{}}{\scriptsize Note: In answer-level
rows, $q=q(a)$. Standard errors are clustered at the tag level.
Post is absorbed by the saturated month fixed effects and is not
separately identified in columns that include month fixed effects.
The interaction terms remain identified as post-period slope
differences relative to the pre-period.
$^{*}p<0.05$, $^{**}p<0.01$, $^{***}p<0.001$.}
\end{tabular}
\end{table*}

% Online appendix / supplementary tables for IJIMDI submission.
% Extracted verbatim from the original manuscript; empirical values unchanged.
% Included from the main manuscript via % Additional tables relocated from the main manuscript for length.
\section{Additional main-text tables}
\label{app:main_tables}

Table~\ref{tab:descriptives} reports answer-level descriptive statistics.
Table~\ref{tab:baseline} reports baseline Post$\times$ASI associations.
Table~\ref{tab:ftests} reports joint event-study $F$-tests.
Table~\ref{tab:hsummary} summarizes primary and secondary tests.

\begin{table*}[t]
\centering
\caption{Descriptive statistics (answer-level sample, $N = 4{,}215{,}082$)}
\label{tab:descriptives}
\scriptsize
\setlength{\tabcolsep}{2pt}
\begin{tabular}{lcccccccc}
\toprule
Variable & $N$ & Mean & SD & P25 & Median & P75 & Min & Max \\
\midrule
Net vote score & 4,215,082 & 3.14 & 8.42 & 0 & 1 & 4 & $-12$ & 1,405 \\
Accepted answer (binary) & 4,215,082 & 0.42 & 0.49 & 0 & 0 & 1 & 0 & 1 \\
Answer comments & 4,215,082 & 1.85 & 3.21 & 0 & 1 & 2 & 0 & 84 \\
Informational index ($\mathrm{ICI}_a$) & 4,215,082 & 0.64 & 0.22 & 0.48 & 0.65 & 0.82 & 0 & 1 \\
Epistemic index ($\mathrm{ECI}_a$) & 4,215,082 & 0.38 & 0.25 & 0.18 & 0.34 & 0.55 & 0 & 1 \\
AI answerability ($\mathrm{ASI}_q$) & 1,842,110 & 0.55 & 0.28 & 0.32 & 0.56 & 0.78 & 0 & 1 \\
Complexity & 1,842,110 & 0.41 & 0.30 & 0.15 & 0.38 & 0.62 & 0 & 1 \\
Pre-event reputation (raw source variable) & 354,209 & 4,520 & 12,450 & 15 & 410 & 3,850 & 1 & 1.2M \\
Tag volatility & 14,502 & 0.22 & 0.18 & 0.08 & 0.19 & 0.31 & 0 & 0.88 \\
\bottomrule
\end{tabular}
\end{table*}

\begin{table}[t]
\centering
\caption{Joint $F$-tests for event-study dynamics and pre-trends}
\label{tab:ftests}
\scriptsize
\setlength{\tabcolsep}{4pt}
\begin{tabular}{llccc}
\toprule
Test & Coefficients tested & Num.\ df & $F$ & $p$ \\
\midrule
Local pre-trend leads & $k=-6$ to $-2$ & 5 & 0.86 & 0.51 \\
Code-LLM leads & $k^{C}=-6$ to $-1$ & 6 & 0.71 & 0.64 \\
Code-LLM lags & $k^{C}=0$ to $+5$ & 6 & 1.27 & 0.27 \\
Eight pre-window bins & Table~\ref{tab:prewindow} & 8 & 1.14 & 0.33 \\
Post-arrival months & $k=0$ to $+12$ & 13 & 4.38 & $<$0.001 \\
\bottomrule
\multicolumn{5}{@{}p{\textwidth}@{}}{\scriptsize Joint $F$-tests evaluate
the indicated event-study bins. Each row names the coefficients in the
contrast. The plotted window and the tested window are not the same.
$p$-values are computed using cluster-robust standard errors.}
\end{tabular}
\end{table}

\begin{table*}[t]
\centering
\caption{Primary tests and bounded extensions}
\label{tab:hsummary}
\scriptsize
\setlength{\tabcolsep}{3pt}
\begin{tabular}{p{0.06\textwidth}p{0.25\textwidth}p{0.18\textwidth}p{0.10\textwidth}p{0.28\textwidth}}
\toprule
Test & Prediction & Estimate (95\% CI) & Raw $p$ & Evidentiary status \\
\midrule
H1 & Informational relative value declines & $-0.153$ $[-0.224,-0.082]$ & 0.00002 & Core; Bonferroni(12) $p=0.00024$ \\
H2 & Epistemic relative value rises & $0.228$ $[0.148,0.308]$ & $<$0.0001 & Core; remains below $0.05$ under Bonferroni(12) \\
H3 & Informational association more negative at high ASI & $-0.139$ $[-0.227,-0.051]$ & 0.0020 & Core; Bonferroni(12) $p=0.024$ \\
H4 & Epistemic association more positive at high ASI & $0.176$ $[0.074,0.278]$ & 0.00071 & Core; Bonferroni(12) $p=0.0086$ \\
H5 & Higher-status users shift toward epistemic work & $0.058$ $[0.034,0.082]$ & $<$0.0001 & Secondary; continuing contributors \\
H6 & Epistemic premium stronger under volatility & $0.079$ $[0.006,0.152]$ & 0.034 & Exploratory; Bonferroni(8) $p=0.272$ \\
H7 & Higher-status users shift toward lower-ASI questions & $-0.082$ $[-0.119,-0.045]$ & $<$0.0001 & Secondary; continuing contributors \\
\bottomrule
\multicolumn{5}{@{}p{\textwidth}@{}}{\scriptsize H1--H4 raw $p$-values are two-sided normal probabilities for the printed coefficient divided by its standard error. H1 is reported as $0.00002$ from $t=-4.25$.}
\end{tabular}
\end{table*}

\begin{table*}[t]
\centering
\caption{Baseline differential post-period associations by AI answerability}
\label{tab:baseline}
\scriptsize
\begin{tabular}{lcccc}
\toprule
& (1) Net vote score (OLS) & (2) Acceptance & (3) Future edits & (4) Survival \\
\midrule
$\mathrm{Post}_t$ & Absorbed & Absorbed & Absorbed & Absorbed \\
$\mathrm{ASI}_q$ & $0.132^{***}$ & $0.032^{*}$ & $0.088^{***}$ & $0.105^{***}$ \\
& (0.033) & (0.014) & (0.021) & (0.028) \\
$\mathrm{Post}_t\times \mathrm{ASI}_q$ & $-0.251^{***}$ & $-0.052^{**}$ & $-0.135^{***}$ & $-0.188^{***}$ \\
& (0.039) & (0.018) & (0.033) & (0.042) \\
Controls/FEs & Yes & Yes & Yes & Yes \\
Observations & 4,215,082 & 4,215,082 & 4,215,082 & 4,215,082 \\
$R^2$ / Pseudo $R^2$ & 0.109 & 0.084 & 0.095 & 0.108 \\
\bottomrule
\multicolumn{5}{@{}p{\textwidth}@{}}{\scriptsize Note: In answer-level
rows, $q=q(a)$. Standard errors are clustered at the tag level.
Post is absorbed by the saturated month fixed effects and is not
separately identified in columns that include month fixed effects.
The interaction terms remain identified as post-period slope
differences relative to the pre-period.
$^{*}p<0.05$, $^{**}p<0.01$, $^{***}p<0.001$.}
\end{tabular}
\end{table*}

% Online appendix / supplementary tables for IJIMDI submission.
% Extracted verbatim from the original manuscript; empirical values unchanged.
% Included from the main manuscript via % Additional tables relocated from the main manuscript for length.
\section{Additional main-text tables}
\label{app:main_tables}

Table~\ref{tab:descriptives} reports answer-level descriptive statistics.
Table~\ref{tab:baseline} reports baseline Post$\times$ASI associations.
Table~\ref{tab:ftests} reports joint event-study $F$-tests.
Table~\ref{tab:hsummary} summarizes primary and secondary tests.

\begin{table*}[t]
\centering
\caption{Descriptive statistics (answer-level sample, $N = 4{,}215{,}082$)}
\label{tab:descriptives}
\scriptsize
\setlength{\tabcolsep}{2pt}
\begin{tabular}{lcccccccc}
\toprule
Variable & $N$ & Mean & SD & P25 & Median & P75 & Min & Max \\
\midrule
Net vote score & 4,215,082 & 3.14 & 8.42 & 0 & 1 & 4 & $-12$ & 1,405 \\
Accepted answer (binary) & 4,215,082 & 0.42 & 0.49 & 0 & 0 & 1 & 0 & 1 \\
Answer comments & 4,215,082 & 1.85 & 3.21 & 0 & 1 & 2 & 0 & 84 \\
Informational index ($\mathrm{ICI}_a$) & 4,215,082 & 0.64 & 0.22 & 0.48 & 0.65 & 0.82 & 0 & 1 \\
Epistemic index ($\mathrm{ECI}_a$) & 4,215,082 & 0.38 & 0.25 & 0.18 & 0.34 & 0.55 & 0 & 1 \\
AI answerability ($\mathrm{ASI}_q$) & 1,842,110 & 0.55 & 0.28 & 0.32 & 0.56 & 0.78 & 0 & 1 \\
Complexity & 1,842,110 & 0.41 & 0.30 & 0.15 & 0.38 & 0.62 & 0 & 1 \\
Pre-event reputation (raw source variable) & 354,209 & 4,520 & 12,450 & 15 & 410 & 3,850 & 1 & 1.2M \\
Tag volatility & 14,502 & 0.22 & 0.18 & 0.08 & 0.19 & 0.31 & 0 & 0.88 \\
\bottomrule
\end{tabular}
\end{table*}

\begin{table}[t]
\centering
\caption{Joint $F$-tests for event-study dynamics and pre-trends}
\label{tab:ftests}
\scriptsize
\setlength{\tabcolsep}{4pt}
\begin{tabular}{llccc}
\toprule
Test & Coefficients tested & Num.\ df & $F$ & $p$ \\
\midrule
Local pre-trend leads & $k=-6$ to $-2$ & 5 & 0.86 & 0.51 \\
Code-LLM leads & $k^{C}=-6$ to $-1$ & 6 & 0.71 & 0.64 \\
Code-LLM lags & $k^{C}=0$ to $+5$ & 6 & 1.27 & 0.27 \\
Eight pre-window bins & Table~\ref{tab:prewindow} & 8 & 1.14 & 0.33 \\
Post-arrival months & $k=0$ to $+12$ & 13 & 4.38 & $<$0.001 \\
\bottomrule
\multicolumn{5}{@{}p{\textwidth}@{}}{\scriptsize Joint $F$-tests evaluate
the indicated event-study bins. Each row names the coefficients in the
contrast. The plotted window and the tested window are not the same.
$p$-values are computed using cluster-robust standard errors.}
\end{tabular}
\end{table}

\begin{table*}[t]
\centering
\caption{Primary tests and bounded extensions}
\label{tab:hsummary}
\scriptsize
\setlength{\tabcolsep}{3pt}
\begin{tabular}{p{0.06\textwidth}p{0.25\textwidth}p{0.18\textwidth}p{0.10\textwidth}p{0.28\textwidth}}
\toprule
Test & Prediction & Estimate (95\% CI) & Raw $p$ & Evidentiary status \\
\midrule
H1 & Informational relative value declines & $-0.153$ $[-0.224,-0.082]$ & 0.00002 & Core; Bonferroni(12) $p=0.00024$ \\
H2 & Epistemic relative value rises & $0.228$ $[0.148,0.308]$ & $<$0.0001 & Core; remains below $0.05$ under Bonferroni(12) \\
H3 & Informational association more negative at high ASI & $-0.139$ $[-0.227,-0.051]$ & 0.0020 & Core; Bonferroni(12) $p=0.024$ \\
H4 & Epistemic association more positive at high ASI & $0.176$ $[0.074,0.278]$ & 0.00071 & Core; Bonferroni(12) $p=0.0086$ \\
H5 & Higher-status users shift toward epistemic work & $0.058$ $[0.034,0.082]$ & $<$0.0001 & Secondary; continuing contributors \\
H6 & Epistemic premium stronger under volatility & $0.079$ $[0.006,0.152]$ & 0.034 & Exploratory; Bonferroni(8) $p=0.272$ \\
H7 & Higher-status users shift toward lower-ASI questions & $-0.082$ $[-0.119,-0.045]$ & $<$0.0001 & Secondary; continuing contributors \\
\bottomrule
\multicolumn{5}{@{}p{\textwidth}@{}}{\scriptsize H1--H4 raw $p$-values are two-sided normal probabilities for the printed coefficient divided by its standard error. H1 is reported as $0.00002$ from $t=-4.25$.}
\end{tabular}
\end{table*}

\begin{table*}[t]
\centering
\caption{Baseline differential post-period associations by AI answerability}
\label{tab:baseline}
\scriptsize
\begin{tabular}{lcccc}
\toprule
& (1) Net vote score (OLS) & (2) Acceptance & (3) Future edits & (4) Survival \\
\midrule
$\mathrm{Post}_t$ & Absorbed & Absorbed & Absorbed & Absorbed \\
$\mathrm{ASI}_q$ & $0.132^{***}$ & $0.032^{*}$ & $0.088^{***}$ & $0.105^{***}$ \\
& (0.033) & (0.014) & (0.021) & (0.028) \\
$\mathrm{Post}_t\times \mathrm{ASI}_q$ & $-0.251^{***}$ & $-0.052^{**}$ & $-0.135^{***}$ & $-0.188^{***}$ \\
& (0.039) & (0.018) & (0.033) & (0.042) \\
Controls/FEs & Yes & Yes & Yes & Yes \\
Observations & 4,215,082 & 4,215,082 & 4,215,082 & 4,215,082 \\
$R^2$ / Pseudo $R^2$ & 0.109 & 0.084 & 0.095 & 0.108 \\
\bottomrule
\multicolumn{5}{@{}p{\textwidth}@{}}{\scriptsize Note: In answer-level
rows, $q=q(a)$. Standard errors are clustered at the tag level.
Post is absorbed by the saturated month fixed effects and is not
separately identified in columns that include month fixed effects.
The interaction terms remain identified as post-period slope
differences relative to the pre-period.
$^{*}p<0.05$, $^{**}p<0.01$, $^{***}p<0.001$.}
\end{tabular}
\end{table*}

% Online appendix / supplementary tables for IJIMDI submission.
% Extracted verbatim from the original manuscript; empirical values unchanged.
% Included from the main manuscript via \input{supplementary/online_appendix}.

\section{Sample construction details}
\label{app:sample}

Table~\ref{tab:sampleledger} records verified analysis-sample counts. Table~\ref{tab:sensitivitysamples} reports H1--H4 coefficients under alternative sample restrictions. Under those alternative exclusions (community-wiki, self-answers, AI-flagged, and all three), the informational coefficient lies between $-0.096$ and $-0.104$ (maximum absolute shift $\Delta\beta_3=0.057$ from baseline $-0.153$). The epistemic shift remains at most $\Delta\beta_4=0.013$. The ``Exclude AI-flagged'' row uses a noisy text-classification proxy (precision 0.74; recall 0.81), not verified authorship.

\textbf{Automated-account rule.} An account was classified as likely automated if it satisfied at least two of three criteria: (i)~more than 50 answers posted within any 24-hour window on three or more separate days; (ii)~pairwise Jaccard similarity $>0.95$ across $\geq$80\% of its answer texts; or (iii)~account creation and first answer within fewer than 90 seconds with no subsequent editing activity. The rule removed 47,218 answers attributed to 3,841 accounts. Manual inspection of 200 randomly sampled flagged accounts confirmed that 194 exhibited bot-like patterns; the remaining 6 were high-volume human contributors. Sensitivity analyses retaining these 6 accounts do not alter the signs or inference of the primary signed-score OLS interactions.

\textbf{Duplicate handling.} Questions linked as duplicates by community closure were retained with a binary duplicate-link indicator included in $\mathbf{X}_{at}$. Exact-text near-duplicates (Jaccard $>0.98$) posted to the \emph{same} question were removed (12,847 records); near-duplicates across different questions were retained.

\begin{table}[t]
\centering
\caption{Verified analysis sample ledger}
\label{tab:sampleledger}
\scriptsize
\begin{tabular}{lr}
\toprule
Field & Verified value \\
\midrule
\texttt{final\_answer\_rows} & 4,215,082 \\
\texttt{unique\_questions} & 1,842,110 \\
\texttt{unique\_contributors} & 354,209 \\
\texttt{post\_answers\_jan23\_dec25} & 1,712,539 \\
\texttt{pre\_and\_transition\_answers} & 2,502,543 \\
\texttt{ai\_flagged\_post\_answers} & 94,217 \\
\texttt{post\_flag\_prevalence} & 5.50\% \\
\texttt{sensitivity\_sample\_exclude\_ai\_flagged} & 4,120,865 \\
\bottomrule
\end{tabular}
\end{table}

\begin{table}[t]
\centering
\caption{Sensitivity of H1--H4 coefficients to sample restrictions}
\label{tab:sensitivitysamples}
\scriptsize
\begin{tabular}{lccc}
\toprule
Restriction & $\beta_3$ (Post$\times$ICI) & $\beta_4$ (Post$\times$ECI) & $N$ \\
\midrule
Baseline (all retained) & $-0.153^{***}$ & $0.228^{***}$ & 4,215,082 \\
Exclude community-wiki & $-0.100^{**}$ & $0.224^{***}$ & 4,127,670 \\
Exclude self-answers & $-0.098^{**}$ & $0.222^{***}$ & 3,791,226 \\
Exclude AI-flagged & $-0.104^{**}$ & $0.231^{***}$ & 4,120,865 \\
Exclude all three & $-0.096^{*}$ & $0.218^{***}$ & 3,703,449 \\
\bottomrule
\multicolumn{4}{@{}p{\textwidth}@{}}{\scriptsize Note: All models include full
controls and fixed effects. ``Exclude AI-flagged'' removes answers selected
by a noisy text-classification proxy (precision 0.74; recall 0.81), not a
verified set of AI-authored answers. The row is therefore a perturbation
test, not evidence that the result holds in an AI-free sample.}
\end{tabular}
\end{table}

\section{AI-assisted text detector}
\label{app:detector}

Likely AI-assisted text was identified using a fine-tuned DeBERTa-v3-large classifier (checkpoint \texttt{microsoft/deberta-v3-large}, Hugging Face transformers 4.38.1) with a binary classification head. The model was trained on 18,000 labeled examples: 9,000 human-written answers sampled from January~2019--October~2022 and 9,000 synthetic answers generated by prompting GPT-3.5-Turbo (March~2023 snapshot) and GPT-4 (August~2023 snapshot) with 4,500 held-out Stack Overflow questions not present in the analytical sample. The positive class is defined as probability $\geq 0.72$, selected by maximizing Youden's $J$ on a 3,000-item validation partition. Training human texts are drawn exclusively from the pre-event window; classification targets are post-January~2023 answers.

Three trained annotators, blinded to detector scores, rated a stratified sample of 1,000 answers. The majority-vote reference is positive for 280 answers and negative for 720. Table~\ref{tab:detectorperf} reports performance. Cell counts are 227 true positives, 53 false negatives, 79 false positives, and 641 true negatives. Performance across sub-periods (2023--2025) showed no statistically significant decline in AUC (range 0.89--0.92; likelihood-ratio test $\chi^2(2)=3.1$, $p=0.21$). Of 1,712,539 post-period answers, 94,217 (5.50\%) exceeded the 0.72 threshold. This figure is an apparent classification rate, not an estimate of true AI-text prevalence. The flag defines sensitivity samples only and is not interpreted as verified authorship.

\begin{table}[t]
\centering
\caption{AI-assisted text detector performance ($n=1{,}000$ human-validation sample)}
\label{tab:detectorperf}
\scriptsize
\begin{tabular}{lc}
\toprule
Metric & Value (95\% CI) \\
\midrule
Sensitivity (recall) & 0.81 [0.76, 0.86] \\
Specificity & 0.89 [0.87, 0.91] \\
Precision & 0.74 [0.69, 0.79] \\
F1-score & 0.77 \\
AUC-ROC & 0.91 [0.89, 0.93] \\
Reference-positive rate & 28\% (280 / 1,000) \\
\bottomrule
\multicolumn{2}{@{}p{\textwidth}@{}}{\scriptsize Counts are 227, 53, 79,
and 641. Sensitivity is $227/280=0.81$, specificity is $641/720=0.89$,
and precision is $227/(227+79)=0.74$. The same precision is
$(0.81\times 0.28)/[(0.81\times 0.28)+(0.11\times 0.72)]=0.74$.
F1 is $2\times 227/(2\times 227+79+53)=0.77$. Wald intervals use those
cell counts.}
\end{tabular}
\end{table}

\section{Notation and coding anchors}
\label{app:notation}

Table~\ref{tab:notation} records principal indices and model notation. Table~\ref{tab:codinganchors} reports coding anchors for informational and epistemic contribution.

\begin{table}[t]
\centering
\caption{Principal indices and model notation}
\label{tab:notation}
\scriptsize
\begin{tabular}{@{}p{0.16\textwidth}p{0.46\textwidth}p{0.26\textwidth}@{}}
\toprule
Symbol & Standardized meaning & Primary use \\
\midrule
$a$ & Answer; the focal answer-level observation & H1--H4 \\
$q$; $q(a)$ & Question containing answer $a$ & ASI, complexity \\
$u$; $u(a)$ & Contributor; author of answer $a$ & H5/H7 panel \\
$t$ & Calendar month & Post indicator, FEs \\
$k$ & Event month relative to focal event & Event study \\
$g$ & Contribution function $\in\{\mathrm{I},\mathrm{E},\mathrm{C}\}$ & Theory only \\
$s(a)$ & Fixed-effect stratum: primary-tag $\times$ question-creation year & Eq.~\eqref{eq:did} \\
$Y_{at}$; $Y_{ut}$ & Answer-level; contributor-month outcome & Outcomes \\
$\mathrm{ICI}_a$; $\mathrm{ECI}_a$ & Informational; epistemic index & Valuation \\
$\mathrm{ASI}_q$ & Expert-rated AI answerability of question $q$ & Exposure proxy \\
$\mathrm{Post}_t$ & Post-event month indicator & Change specs \\
$\mathrm{Status}^{\mathrm{pre}}_u$ & Standardized log pre-event reputation & Adaptation \\
$\mathbf{X}_{at}$; $\mathbf{Z}_{ut}$ & Answer-level; contributor-month controls & Covariates \\
$\mathrm{Volatility}_{\operatorname{tag}(q),t}$ & Tag-level volatility & H6 \\
$\mathrm{RV}_{g,t}(\mathrm{ASI}_q)$ & Return to function $g$ & Theory \\
$\beta$; $\beta_k$; $\beta_1$--$\beta_4$; $\rho$ & Parameters & Estimands \\
$\theta_{\mathrm{I}}$; $\theta_{\mathrm{E}}$ & Triple-interaction coefficients & H3, H4 \\
$\boldsymbol{\gamma}$; $\boldsymbol{\kappa}$ & Control coefficient vectors & Controls \\
$\alpha_{s(a)}$; $\delta_t$ & Stratum; month fixed effects & Baseline \\
$\mu_{u(a)}$; $\lambda_{\operatorname{tag}(q(a))}$; $\eta_u$ & Contributor; tag; panel FEs & Valuation/adaptation \\
$\varepsilon_{at}$; $\varepsilon_{ut}$ & Disturbances & Errors \\
$\mathrm{EventTime}_t$; $\mathbb{1}\{\cdot\}$ & Event month; indicator & Event study \\
\bottomrule
\end{tabular}
\end{table}

\begin{table*}[t]
\centering
\caption{Coding anchors for informational and epistemic contribution}
\label{tab:codinganchors}
\scriptsize
\begin{tabular}{p{0.18\textwidth}p{0.34\textwidth}p{0.34\textwidth}}
\toprule
Pattern & Coding anchor & Illustrative answer form \\
\midrule
High ICI, low ECI & Supplies a direct command, code fragment, or standard
procedure without visible tests, assumptions, or limits & ``Use function
X with argument Y.'' \\
Low ICI, high ECI & Diagnoses uncertainty, compares explanations, or
specifies tests and boundary conditions without offering a complete
drop-in solution & ``This depends on the runtime version; test A before
choosing between B and C.'' \\
High ICI, high ECI & Supplies a usable solution and makes its warrant,
scope, and failure conditions inspectable & ``Use X; the example verifies
condition Y, but it fails under Z and requires version V.'' \\
Low ICI, low ECI & Provides neither a usable solution nor inspectable
reasoning & Unsupported assertion, social comment, or restatement of the
question \\
\bottomrule
\end{tabular}
\end{table*}

\section{Identification and robustness tables}
\label{app:identification}

Table~\ref{tab:secular_control} records that a $\mathrm{Post}_t\times\mathrm{PlatformAnswerVolume}_t$ term cannot be separately estimated alongside saturated month fixed effects. Table~\ref{tab:twoevent} reports the two-event piecewise specification within LLM arrival. Table~\ref{tab:prewindow} reports eight pre-window ASI interactions (joint $F(8)=1.14$, $p=0.33$). Table~\ref{tab:timeplacebo} reports placebo-in-time tests. Table~\ref{tab:lowcopilot} reports H1--H4 outside the first code-LLM languages ($N=1{,}475{,}279$). Table~\ref{tab:tagslope} adds a pre-period tag answer-volume slope control.

\begin{table}[t]
\centering
\caption{Platform-wide secular decline and identification}
\label{tab:secular_control}
\scriptsize
\setlength{\tabcolsep}{3pt}
\begin{tabular}{p{0.28\textwidth}p{0.20\textwidth}p{0.20\textwidth}p{0.24\textwidth}}
\toprule
Specification & Post$\times$ICI & Post$\times$ECI & Treatment of secular decline \\
\midrule
Baseline OLS & $-0.153$ & $0.228$ & Saturated calendar-month fixed effects \\
Add Post $\times$ platform-wide monthly answer volume & Not separately identified & Not separately identified & Collinear with calendar-month fixed effects \\
\bottomrule
\end{tabular}
\end{table}

\begin{table}[t]
\centering
\caption{Two-event piecewise specification within LLM arrival}
\label{tab:twoevent}
\scriptsize
\begin{tabular}{p{0.34\textwidth}p{0.28\textwidth}p{0.28\textwidth}}
\toprule
Window & Joint $F$-test ($p$) & Interpretation \\
\midrule
Code-LLM leads ($k^C=-6$ to $-1$) & $F(6)=0.71$ ($p=0.64$) & Joint test does not reject \\
Code-LLM lags ($k^C=0$ to $+5$) & $F(6)=1.27$ ($p=0.27$) & Short window; test does not reject \\
Leads before public LLM arrival ($k^G=-6$ to $-2$) & Same five-lead test as the event study & Not a second test \\
After public LLM arrival ($k^G=0$ to $+12$) & $F(13)=4.38$ ($p<0.001$) & Jointly significant association \\
\bottomrule
\end{tabular}
\end{table}

\begin{table}[t]
\centering
\caption{Full pre-window $\mathrm{ASI}$ interactions
(net vote score)}
\label{tab:prewindow}
\scriptsize
\begin{tabular}{lcc}
\toprule
Bin, relative to November 2022 & $\beta$ & SE \\
\midrule
January--June 2019 & 0.014 & 0.046 \\
July--December 2019 & $-0.021$ & 0.041 \\
January--June 2020 & 0.008 & 0.037 \\
July--December 2020 & $-0.033$ & 0.034 \\
January--June 2021 & 0.011 & 0.032 \\
July--December 2021 & $-0.027$ & 0.031 \\
January--May 2022 & $-0.019$ & 0.033 \\
June--October 2022 & $-0.046$ & 0.035 \\
\midrule
Joint test of the eight bins & \multicolumn{2}{c}{$F(8)=1.14$, $p=0.33$} \\
\bottomrule
\multicolumn{3}{@{}p{\textwidth}@{}}{\scriptsize The June--October 2022 bin
overlaps early code-LLM diffusion and is the least quiet pre-arrival bin.}
\end{tabular}
\end{table}

\begin{table}[t]
\centering
\caption{Placebo-in-time tests, pre-arrival sample only}
\label{tab:timeplacebo}
\scriptsize
\begin{tabular}{lcccc}
\toprule
Pseudo-event & Post$\times$ICI & SE & Post$\times$ECI & SE \\
\midrule
30 November 2019 & $-0.016$ & 0.047 & 0.024 & 0.061 \\
30 November 2020 & $-0.028$ & 0.043 & 0.019 & 0.056 \\
30 June 2021 & $-0.034$ & 0.041 & 0.038 & 0.053 \\
\midrule
LLM arrival (30 November 2022) & $-0.153$ & 0.036 & 0.228 & 0.041 \\
\bottomrule
\multicolumn{5}{@{}p{\textwidth}@{}}{\scriptsize The last row is the
reported arrival estimate. Each placebo coefficient is smaller than twice
its printed standard error. The largest absolute ratio of coefficient to
standard error is $0.83$. No stored $p$-value is reported for these
cells. The printed standard errors do not indicate a departure at $0.05$,
and the comparison does not prove that the interactions are unique to
late 2022.}
\end{tabular}
\end{table}

\begin{table}[t]
\centering
\caption{H1--H4 estimates outside the first code-LLM languages}
\label{tab:lowcopilot}
\scriptsize
\begin{tabular}{lccc}
\toprule
Coefficient & Reported, full sample & Outside first code-LLM languages & SE \\
\midrule
Post$\times$ICI & $-0.153$ & $-0.086$ & 0.049 \\
Post$\times$ECI & $0.228$ & $0.191$ & 0.064 \\
Post$\times$ASI$\times$ICI & $-0.139$ & $-0.117$ & 0.055 \\
Post$\times$ASI$\times$ECI & $0.176$ & $0.154$ & 0.058 \\
\bottomrule
\multicolumn{4}{@{}p{\textwidth}@{}}{\scriptsize Subsample
$N=1{,}475{,}279$.}
\end{tabular}
\end{table}

\begin{table}[t]
\centering
\caption{Control for pre-period tag answer-volume slopes}
\label{tab:tagslope}
\scriptsize
\begin{tabular}{lcc}
\toprule
Coefficient & Reported & With tag slope \\
\midrule
Post$\times$ICI & $-0.153$ (0.036) & $-0.087$ (0.039) \\
Post$\times$ECI & $0.228$ (0.041) & $0.194$ (0.053) \\
Post$\times$ASI$\times$ICI & $-0.139$ (0.045) & $-0.121$ (0.044) \\
Post$\times$ASI$\times$ECI & $0.176$ (0.052) & $0.149$ (0.046) \\
\bottomrule
\multicolumn{3}{@{}p{\textwidth}@{}}{\scriptsize Parentheses are standard
errors.}
\end{tabular}
\end{table}

\section{Inference, diagnostics, and multiplicity}
\label{app:inference}

Table~\ref{tab:pvalues} reports four-test and eight-test Bonferroni values for the primary net-vote family. Table~\ref{tab:family12} reports the twelve-test Bonferroni family. Table~\ref{tab:diagnostics} summarizes primary model diagnostics. Table~\ref{tab:vif} reports generalized VIFs.

\begin{table}[t]
\centering
\caption{Primary net-vote $p$-values with retrospective multiplicity corrections}
\label{tab:pvalues}
\scriptsize
\begin{tabular}{lccc}
\toprule
H & Raw $p$ & Bonferroni, four tests & Bonferroni, eight-test sensitivity \\
\midrule
H1 & 0.00002 & 0.00008 & 0.00016 \\
H2 & $<$0.0001 & $<$0.0004 & $<$0.0008 \\
H3 & 0.0020 & 0.0080 & 0.016 \\
H4 & 0.00071 & 0.0029 & 0.0057 \\
\bottomrule
\end{tabular}
\end{table}

\begin{table}[t]
\centering
\caption{Twelve-test Bonferroni values computed from printed estimates}
\label{tab:family12}
\scriptsize
\begin{tabular}{llcc}
\toprule
Coefficient & Source $p$ & $\times 12$ & Below 0.05 \\
\midrule
Post$\times$ICI, net votes & 0.00002 & 0.00024 & Yes \\
Post$\times$ECI, net votes & $<$0.0001 & $<$0.0012 & Yes \\
Post$\times$ASI$\times$ICI, net votes & 0.0020 & 0.024 & Yes \\
Post$\times$ASI$\times$ECI, net votes & 0.00071 & 0.0086 & Yes \\
Post$\times$ICI, acceptance & 0.031 & 0.375 & No \\
Post$\times$ECI, acceptance & 0.0003 & 0.004 & Yes \\
Post$\times$ASI$\times$ICI, acceptance & 0.006 & 0.075 & No \\
Post$\times$ASI$\times$ECI, acceptance & 0.0006 & 0.008 & Yes \\
H5, ECI share & $<$0.0001 & $<$0.0012 & Yes \\
H7, selected ASI & $<$0.0001 & $<$0.0012 & Yes \\
H6, volatility & 0.034 & 0.408 & No \\
Post$\times$status, count & 0.0014 & 0.017 & Yes \\
\bottomrule
\multicolumn{4}{@{}p{\textwidth}@{}}{\scriptsize Acceptance and count
$p$-values are two-sided Wald tests computed from the reported
coefficients and standard errors ($z=\hat\beta/\mathrm{SE}$).
H1, H3, and H4 net-vote $p$-values are two-sided normal probabilities
for the printed coefficient divided by its standard error.
H1 uses $t=-0.153/0.036=-4.25$, reported as $0.00002$; twelve times
that reported value is $0.00024$. H3 uses $t=-0.139/0.045=-3.09$
($p=0.0020$). H4 uses $t=0.176/0.052=3.38$ ($p=0.00071$).}
\end{tabular}
\end{table}

\begin{table*}[t]
\centering
\caption{Primary model diagnostics}
\label{tab:diagnostics}
\scriptsize
\setlength{\tabcolsep}{3pt}
\begin{tabular}{p{0.22\textwidth}p{0.17\textwidth}p{0.17\textwidth}p{0.17\textwidth}p{0.17\textwidth}}
\toprule
Diagnostic & Net vote score (OLS) & Acceptance (logit) & ECI share (frac.\ resp.) & Exit hazard (cloglog) \\
\midrule
Estimator & OLS & Logit & Fractional response & Complementary log--log \\
Clustering level & Tag (14,502) & Contributor (354,209) & Contributor & Contributor \\
Outcome mean / variance & 3.14 / 70.9 & 0.42 / 0.24 & 0.38 / 0.06 & Event rate 3.2\% \\
Distribution / zero share & Median 1; zero share 38.2\%; min $-12$, max 1,405 & 0 to 1 & 0 to 1 & 0 to 1 \\
Robust variance estimator & Tag-clustered HC & Contributor-clustered & Contributor-clustered & Contributor-clustered \\
Heteroskedasticity & White test $p<0.001$; tag-clustered robust SE & --- & --- & --- \\
Convergence & Closed-form OLS solution & No perfect prediction; Firth check stable & Converged & Converged \\
Link/specification test & RESET-type $p=0.26$ & Hosmer--Lemeshow $p=0.41$ & RESET-type $p=0.31$ & Time-varying coeff.\ $p=0.34$ \\
Influence / cluster sensitivity & Leave-one-tag-out maximum coefficient change $=0.019$ & --- & --- & --- \\
Calibration (Brier; slope) & --- & 0.21; 0.97 & --- & --- \\
Discrimination & --- & AUC = 0.71 & --- & C-index = 0.74 (attrition) \\
Omnibus Wald test & $p<0.001$ & $p<0.001$ & $p<0.001$ & $p<0.001$ \\
\bottomrule
\end{tabular}
\end{table*}

\begin{table}[t]
\centering
\caption{Generalized VIFs for the valuation specification and the separate supply specification (centered)}
\label{tab:vif}
\scriptsize
\begin{tabular}{lccc}
\toprule
Term & Df & GVIF & GVIF$^{1/(2\,\mathrm{Df})}$ \\
\midrule
\multicolumn{4}{l}{\textit{Valuation specification}} \\
$\mathrm{Post}_t$ & 1 & 1.02 & 1.01 \\
$\mathrm{ICI}_a$ & 1 & 1.18 & 1.09 \\
$\mathrm{ECI}_a$ & 1 & 1.24 & 1.11 \\
$\mathrm{ASI}_{q(a)}$ & 1 & 1.15 & 1.07 \\
$\mathrm{Post}_t\times\mathrm{ICI}_a$ & 1 & 2.31 & 1.52 \\
$\mathrm{Post}_t\times\mathrm{ECI}_a$ & 1 & 2.48 & 1.57 \\
$\mathrm{Post}_t\times\mathrm{ASI}_{q(a)}$ & 1 & 2.12 & 1.46 \\
$\mathrm{Post}_t\times\mathrm{ASI}_{q(a)}\times\mathrm{ICI}_a$ & 1 & 3.85 & 1.96 \\
$\mathrm{Post}_t\times\mathrm{ASI}_{q(a)}\times\mathrm{ECI}_a$ & 1 & 3.92 & 1.98 \\
\midrule
\multicolumn{4}{l}{\textit{Supply specification}} \\
ICI$\times$Post$\times$Supply & 1 & 3.84 & 1.96 \\
ECI$\times$Post$\times$Supply & 1 & 5.29 & 2.30 \\
\bottomrule
\end{tabular}
\par\smallskip
{\scriptsize Note: Df is the number of model-matrix columns assigned to
the term. Each displayed continuous interaction contributes one
column, so the adjusted value is the square root of the reported GVIF.
The two panels are separate specifications.}
\end{table}

\section{Measurement and ASI validation}
\label{app:measurement_ext}

Table~\ref{tab:modelasi} reports agreement between expert ASI and a frozen pre-30~November~2022 model checkpoint on the 600-question expert panel. Correctness rises from $0.19$ in the lowest ASI decile to $0.67$ in the highest; ASI correlates $0.58$ with correctness and $0.51$ with usability. Estimating H3 and H4 on the $N=1{,}380$ attached answers yields $\mathrm{Post}\times\mathrm{Correct}\times\mathrm{ICI}=-0.11$ (SE $0.09$) and $\mathrm{Post}\times\mathrm{Correct}\times\mathrm{ECI}=0.15$ (SE $0.10$), same sign as $-0.139$ and $0.176$ but too imprecise to confirm the answerability interpretation alone. Raising the AI-flag threshold to precision $0.85$ yields Post$\times$ICI $=-0.100$ (SE $0.037$) and Post$\times$ECI $=0.224$ (SE $0.050$), with $N=4{,}153{,}400$.

\begin{table}[t]
\centering
\caption{Agreement between expert ASI and a frozen model
($n=600$ questions)}
\label{tab:modelasi}
\scriptsize
\begin{tabular}{lc}
\toprule
Comparison & Estimate \\
\midrule
ASI with model correctness ($r$) & 0.58 \\
ASI with model usability ($r$) & 0.51 \\
ASI rank with correctness (Spearman) & 0.54 \\
Correctness, ASI decile 1 / decile 10 & 0.19 / 0.67 \\
Usability, ASI decile 1 / decile 10 & 0.24 / 0.61 \\
\bottomrule
\multicolumn{2}{@{}p{\textwidth}@{}}{\scriptsize The 600-question sample is
the expert panel used for ASI validation.}
\end{tabular}
\end{table}

\section{Additional results and sensitivity analyses}
\label{app:extra_results}

Table~\ref{tab:outcomerobust} reports top-1\% trim and ordered outcome specifications. Table~\ref{tab:attrition} and Table~\ref{tab:ipw} document attrition and IPW adaptation estimates. Table~\ref{tab:placebo} reports a placebo outcome. Table~\ref{tab:supply} reports an exploratory informational-supply association (not causal mediation). Table~\ref{tab:exitbounds} reports person-weighted exit scenarios. Table~\ref{tab:deleted} reports hypothetical deleted-answer bounds; these are not estimated regression coefficients and cannot recover deleted posts.

\begin{table*}[t]
\centering
\caption{Robustness of the primary interactions to alternative outcome specifications}
\label{tab:outcomerobust}
\scriptsize
\setlength{\tabcolsep}{4pt}
\begin{tabular}{llccc}
\toprule
Hypothesis & Variable & Baseline (OLS) & Trim top 1\% (OLS) & Ordered outcome (logit) \\
\midrule
H1 & Post$\times$ICI & $-0.153$ (0.036)$^{***}$ & $-0.141$ (0.038)$^{***}$ & $-0.412$ (0.098)$^{***}$ \\
H2 & Post$\times$ECI & $0.228$ (0.041)$^{***}$ & $0.215$ (0.043)$^{***}$ & $0.615$ (0.110)$^{***}$ \\
H3 & Post$\times$ICI$\times$ASI & $-0.139$ (0.045)$^{**}$ & $-0.125$ (0.048)$^{**}$ & $-0.380$ (0.125)$^{**}$ \\
H4 & Post$\times$ECI$\times$ASI & $0.176$ (0.052)$^{***}$ & $0.162$ (0.055)$^{**}$ & $0.475$ (0.140)$^{***}$ \\
\midrule
& Observations & 4,215,082 & 4,172,931 & 4,215,082 \\
& Adjusted $R^2$ / pseudo $R^2$ & 0.142 & 0.138 & 0.095 \\
\bottomrule
\multicolumn{5}{@{}p{\textwidth}@{}}{\scriptsize Standard errors in
parentheses are clustered at the tag level. $^{***}p<0.001$,
$^{**}p<0.01$, from the normal tail of the printed coefficient divided
by its standard error. Baseline cells repeat the net-vote column of
Table~\ref{tab:epistemic}. The trim drops the top 1\% of net-vote
answers. The ordered outcome uses a five-category vote binning. Ordered
coefficients are on the latent index.}
\end{tabular}
\end{table*}

\begin{table}[t]
\centering
\caption{Pre-event characteristics of retained versus exiting contributors}
\label{tab:attrition}
\scriptsize
\begin{tabular}{lcc}
\toprule
Pre-event characteristic & Retained ($n=218{,}447$) & Exited ($n=135{,}762$) \\
\midrule
Median pre-event reputation & 410 & 185 \\
Mean monthly answers (pre) & 2.8 & 1.1 \\
Mean ECI share & 0.38 & 0.31 \\
Mean selected ASI & 0.55 & 0.61 \\
Mean tag entropy & 1.9 & 1.2 \\
\bottomrule
\end{tabular}
\end{table}

\begin{table}[t]
\centering
\caption{Attrition-adjusted adaptation estimates}
\label{tab:ipw}
\scriptsize
\begin{tabular}{lccc}
\toprule
Outcome & Unweighted & IPW & Cohort (explicit zeros) \\
\midrule
ECI share & $0.058^{***}$ & $0.051^{***}$ & $0.047^{**}$ \\
& (0.012) & (0.014) & (0.015) \\
Selected ASI & $-0.082^{***}$ & $-0.074^{***}$ & $-0.069^{**}$ \\
& (0.019) & (0.022) & (0.024) \\
Complexity & $0.045^{**}$ & $0.039^{*}$ & $0.036^{*}$ \\
& (0.015) & (0.017) & (0.018) \\
\bottomrule
\multicolumn{4}{@{}p{\textwidth}@{}}{\scriptsize $^{*}p<0.05$, $^{**}p<0.01$, $^{***}p<0.001$.
The cohort epistemic-share estimate is $0.047/0.015=3.13$ (two-sided $p=0.0017$), so it takes two stars.
The cohort specification retains all 354,209 pre-period contributors.}
\end{tabular}
\end{table}

\begin{table}[t]
\centering
\caption{Placebo test: pre-period edit volume}
\label{tab:placebo}
\scriptsize
\begin{tabular}{lccc}
\toprule
Outcome & Interaction & Estimate & $p$-value \\
\midrule
Pre-period answer edit volume & Post$\times$ASI & 0.004 & 0.612 \\
\bottomrule
\multicolumn{4}{@{}p{\textwidth}@{}}{\scriptsize Note: The coefficient is
interpreted as a placebo differential association, not as proof of a
structural zero.}
\end{tabular}
\end{table}

\begin{table}[t]
\centering
\caption{Informational-supply association}
\label{tab:supply}
\scriptsize
\begin{tabular}{lc}
\toprule
Quantity & Estimate \\
\midrule
Post-period change in supply, high-ASI tags (log points) & 0.17 \\
Post-period change in supply, low-ASI tags (log points) & 0.04 \\
Supply$\times$ECI & 0.092 (SE 0.040) \\
Post$\times$ICI after supply control & $-0.077$ (SE 0.038) \\
Post$\times$ECI after supply control & 0.171 (SE 0.052) \\
Post$\times$ASI$\times$ECI after supply control & 0.128 (SE 0.046) \\
Informational interquartile contrast after supply control & $-0.026$ \\
Share of the $-0.153$ coefficient remaining & 0.503 \\
Share of the $0.228$ coefficient remaining & 0.75 \\
\bottomrule
\multicolumn{2}{@{}p{\textwidth}@{}}{\scriptsize High- and low-ASI tags are
the top and bottom quartiles of mean pre-period ASI. The supply term is
not interpreted as a causal mediator.}
\end{tabular}
\end{table}

\begin{table}[t]
\centering
\caption{Person-weighted scenarios for the status interactions}
\label{tab:exitbounds}
\scriptsize
\begin{tabular}{lcc}
\toprule
Assumption for the 135,762 exiters & H5, ECI share & H7, selected ASI \\
\midrule
Same shift as stayers (reported $0.058$, $-0.082$) & 0.058 & $-0.082$ \\
Half the stayer shift & 0.047 & $-0.066$ \\
No shift & 0.036 & $-0.051$ \\
Opposite shift of the same size & 0.014 & $-0.019$ \\
\bottomrule
\multicolumn{3}{@{}p{\textwidth}@{}}{\scriptsize Weights are $0.617$ on
stayers and $0.383$ on exiters. Calculated from the reported
coefficients and the attrition counts. Not an estimate of exiter
behavior.}
\end{tabular}
\end{table}

\begin{table*}[t]
\centering
\caption{Hypothetical sensitivity bounds under assumed deleted-answer scenarios}
\label{tab:deleted}
\scriptsize
\setlength{\tabcolsep}{3pt}
\begin{tabular}{lcccccc}
\toprule
Scenario & \makecell{Post-period\\deletion} & \makecell{Pre-period\\deletion} & \makecell{Vote penalty\\high ICI} & \makecell{Vote penalty\\high ECI} & \makecell{Hypothetical\\Post$\times$ICI bound} & \makecell{Hypothetical\\Post$\times$ECI bound} \\
\midrule
1 Baseline & 8\% & 3\% & $-1.8$ & $-0.4$ & $-0.119$ & $0.211$ \\
2 Conservative & 5\% & 2\% & $-1.8$ & $-0.4$ & $-0.131$ & $0.218$ \\
3 Aggressive & 10\% & 5\% & $-1.8$ & $-0.4$ & $-0.105$ & $0.202$ \\
4 High penalty & 8\% & 3\% & $-2.5$ & $-0.6$ & $-0.108$ & $0.195$ \\
5 Worst case & 10\% & 5\% & $-2.5$ & $-0.6$ & $-0.092$ & $0.184$ \\
6 Minimal impact & 2\% & 1\% & $-1.0$ & $-0.2$ & $-0.142$ & $0.225$ \\
\bottomrule
\multicolumn{7}{@{}p{\textwidth}@{}}{\scriptsize These values are
hypothetical bounds from applying assumed vote penalties to assumed
deletion rates. They are not estimated regression coefficients, and they
cannot be reproduced from the exported covariance matrices. The signs
match the reported $-0.153$ and $0.228$ in every row. The bounds
illustrate directional robustness under those assumptions. They do not
recover the deleted posts.}
\end{tabular}
\end{table*}
.

\section{Sample construction details}
\label{app:sample}

Table~\ref{tab:sampleledger} records verified analysis-sample counts. Table~\ref{tab:sensitivitysamples} reports H1--H4 coefficients under alternative sample restrictions. Under those alternative exclusions (community-wiki, self-answers, AI-flagged, and all three), the informational coefficient lies between $-0.096$ and $-0.104$ (maximum absolute shift $\Delta\beta_3=0.057$ from baseline $-0.153$). The epistemic shift remains at most $\Delta\beta_4=0.013$. The ``Exclude AI-flagged'' row uses a noisy text-classification proxy (precision 0.74; recall 0.81), not verified authorship.

\textbf{Automated-account rule.} An account was classified as likely automated if it satisfied at least two of three criteria: (i)~more than 50 answers posted within any 24-hour window on three or more separate days; (ii)~pairwise Jaccard similarity $>0.95$ across $\geq$80\% of its answer texts; or (iii)~account creation and first answer within fewer than 90 seconds with no subsequent editing activity. The rule removed 47,218 answers attributed to 3,841 accounts. Manual inspection of 200 randomly sampled flagged accounts confirmed that 194 exhibited bot-like patterns; the remaining 6 were high-volume human contributors. Sensitivity analyses retaining these 6 accounts do not alter the signs or inference of the primary signed-score OLS interactions.

\textbf{Duplicate handling.} Questions linked as duplicates by community closure were retained with a binary duplicate-link indicator included in $\mathbf{X}_{at}$. Exact-text near-duplicates (Jaccard $>0.98$) posted to the \emph{same} question were removed (12,847 records); near-duplicates across different questions were retained.

\begin{table}[t]
\centering
\caption{Verified analysis sample ledger}
\label{tab:sampleledger}
\scriptsize
\begin{tabular}{lr}
\toprule
Field & Verified value \\
\midrule
\texttt{final\_answer\_rows} & 4,215,082 \\
\texttt{unique\_questions} & 1,842,110 \\
\texttt{unique\_contributors} & 354,209 \\
\texttt{post\_answers\_jan23\_dec25} & 1,712,539 \\
\texttt{pre\_and\_transition\_answers} & 2,502,543 \\
\texttt{ai\_flagged\_post\_answers} & 94,217 \\
\texttt{post\_flag\_prevalence} & 5.50\% \\
\texttt{sensitivity\_sample\_exclude\_ai\_flagged} & 4,120,865 \\
\bottomrule
\end{tabular}
\end{table}

\begin{table}[t]
\centering
\caption{Sensitivity of H1--H4 coefficients to sample restrictions}
\label{tab:sensitivitysamples}
\scriptsize
\begin{tabular}{lccc}
\toprule
Restriction & $\beta_3$ (Post$\times$ICI) & $\beta_4$ (Post$\times$ECI) & $N$ \\
\midrule
Baseline (all retained) & $-0.153^{***}$ & $0.228^{***}$ & 4,215,082 \\
Exclude community-wiki & $-0.100^{**}$ & $0.224^{***}$ & 4,127,670 \\
Exclude self-answers & $-0.098^{**}$ & $0.222^{***}$ & 3,791,226 \\
Exclude AI-flagged & $-0.104^{**}$ & $0.231^{***}$ & 4,120,865 \\
Exclude all three & $-0.096^{*}$ & $0.218^{***}$ & 3,703,449 \\
\bottomrule
\multicolumn{4}{@{}p{\textwidth}@{}}{\scriptsize Note: All models include full
controls and fixed effects. ``Exclude AI-flagged'' removes answers selected
by a noisy text-classification proxy (precision 0.74; recall 0.81), not a
verified set of AI-authored answers. The row is therefore a perturbation
test, not evidence that the result holds in an AI-free sample.}
\end{tabular}
\end{table}

\section{AI-assisted text detector}
\label{app:detector}

Likely AI-assisted text was identified using a fine-tuned DeBERTa-v3-large classifier (checkpoint \texttt{microsoft/deberta-v3-large}, Hugging Face transformers 4.38.1) with a binary classification head. The model was trained on 18,000 labeled examples: 9,000 human-written answers sampled from January~2019--October~2022 and 9,000 synthetic answers generated by prompting GPT-3.5-Turbo (March~2023 snapshot) and GPT-4 (August~2023 snapshot) with 4,500 held-out Stack Overflow questions not present in the analytical sample. The positive class is defined as probability $\geq 0.72$, selected by maximizing Youden's $J$ on a 3,000-item validation partition. Training human texts are drawn exclusively from the pre-event window; classification targets are post-January~2023 answers.

Three trained annotators, blinded to detector scores, rated a stratified sample of 1,000 answers. The majority-vote reference is positive for 280 answers and negative for 720. Table~\ref{tab:detectorperf} reports performance. Cell counts are 227 true positives, 53 false negatives, 79 false positives, and 641 true negatives. Performance across sub-periods (2023--2025) showed no statistically significant decline in AUC (range 0.89--0.92; likelihood-ratio test $\chi^2(2)=3.1$, $p=0.21$). Of 1,712,539 post-period answers, 94,217 (5.50\%) exceeded the 0.72 threshold. This figure is an apparent classification rate, not an estimate of true AI-text prevalence. The flag defines sensitivity samples only and is not interpreted as verified authorship.

\begin{table}[t]
\centering
\caption{AI-assisted text detector performance ($n=1{,}000$ human-validation sample)}
\label{tab:detectorperf}
\scriptsize
\begin{tabular}{lc}
\toprule
Metric & Value (95\% CI) \\
\midrule
Sensitivity (recall) & 0.81 [0.76, 0.86] \\
Specificity & 0.89 [0.87, 0.91] \\
Precision & 0.74 [0.69, 0.79] \\
F1-score & 0.77 \\
AUC-ROC & 0.91 [0.89, 0.93] \\
Reference-positive rate & 28\% (280 / 1,000) \\
\bottomrule
\multicolumn{2}{@{}p{\textwidth}@{}}{\scriptsize Counts are 227, 53, 79,
and 641. Sensitivity is $227/280=0.81$, specificity is $641/720=0.89$,
and precision is $227/(227+79)=0.74$. The same precision is
$(0.81\times 0.28)/[(0.81\times 0.28)+(0.11\times 0.72)]=0.74$.
F1 is $2\times 227/(2\times 227+79+53)=0.77$. Wald intervals use those
cell counts.}
\end{tabular}
\end{table}

\section{Notation and coding anchors}
\label{app:notation}

Table~\ref{tab:notation} records principal indices and model notation. Table~\ref{tab:codinganchors} reports coding anchors for informational and epistemic contribution.

\begin{table}[t]
\centering
\caption{Principal indices and model notation}
\label{tab:notation}
\scriptsize
\begin{tabular}{@{}p{0.16\textwidth}p{0.46\textwidth}p{0.26\textwidth}@{}}
\toprule
Symbol & Standardized meaning & Primary use \\
\midrule
$a$ & Answer; the focal answer-level observation & H1--H4 \\
$q$; $q(a)$ & Question containing answer $a$ & ASI, complexity \\
$u$; $u(a)$ & Contributor; author of answer $a$ & H5/H7 panel \\
$t$ & Calendar month & Post indicator, FEs \\
$k$ & Event month relative to focal event & Event study \\
$g$ & Contribution function $\in\{\mathrm{I},\mathrm{E},\mathrm{C}\}$ & Theory only \\
$s(a)$ & Fixed-effect stratum: primary-tag $\times$ question-creation year & Eq.~\eqref{eq:did} \\
$Y_{at}$; $Y_{ut}$ & Answer-level; contributor-month outcome & Outcomes \\
$\mathrm{ICI}_a$; $\mathrm{ECI}_a$ & Informational; epistemic index & Valuation \\
$\mathrm{ASI}_q$ & Expert-rated AI answerability of question $q$ & Exposure proxy \\
$\mathrm{Post}_t$ & Post-event month indicator & Change specs \\
$\mathrm{Status}^{\mathrm{pre}}_u$ & Standardized log pre-event reputation & Adaptation \\
$\mathbf{X}_{at}$; $\mathbf{Z}_{ut}$ & Answer-level; contributor-month controls & Covariates \\
$\mathrm{Volatility}_{\operatorname{tag}(q),t}$ & Tag-level volatility & H6 \\
$\mathrm{RV}_{g,t}(\mathrm{ASI}_q)$ & Return to function $g$ & Theory \\
$\beta$; $\beta_k$; $\beta_1$--$\beta_4$; $\rho$ & Parameters & Estimands \\
$\theta_{\mathrm{I}}$; $\theta_{\mathrm{E}}$ & Triple-interaction coefficients & H3, H4 \\
$\boldsymbol{\gamma}$; $\boldsymbol{\kappa}$ & Control coefficient vectors & Controls \\
$\alpha_{s(a)}$; $\delta_t$ & Stratum; month fixed effects & Baseline \\
$\mu_{u(a)}$; $\lambda_{\operatorname{tag}(q(a))}$; $\eta_u$ & Contributor; tag; panel FEs & Valuation/adaptation \\
$\varepsilon_{at}$; $\varepsilon_{ut}$ & Disturbances & Errors \\
$\mathrm{EventTime}_t$; $\mathbb{1}\{\cdot\}$ & Event month; indicator & Event study \\
\bottomrule
\end{tabular}
\end{table}

\begin{table*}[t]
\centering
\caption{Coding anchors for informational and epistemic contribution}
\label{tab:codinganchors}
\scriptsize
\begin{tabular}{p{0.18\textwidth}p{0.34\textwidth}p{0.34\textwidth}}
\toprule
Pattern & Coding anchor & Illustrative answer form \\
\midrule
High ICI, low ECI & Supplies a direct command, code fragment, or standard
procedure without visible tests, assumptions, or limits & ``Use function
X with argument Y.'' \\
Low ICI, high ECI & Diagnoses uncertainty, compares explanations, or
specifies tests and boundary conditions without offering a complete
drop-in solution & ``This depends on the runtime version; test A before
choosing between B and C.'' \\
High ICI, high ECI & Supplies a usable solution and makes its warrant,
scope, and failure conditions inspectable & ``Use X; the example verifies
condition Y, but it fails under Z and requires version V.'' \\
Low ICI, low ECI & Provides neither a usable solution nor inspectable
reasoning & Unsupported assertion, social comment, or restatement of the
question \\
\bottomrule
\end{tabular}
\end{table*}

\section{Identification and robustness tables}
\label{app:identification}

Table~\ref{tab:secular_control} records that a $\mathrm{Post}_t\times\mathrm{PlatformAnswerVolume}_t$ term cannot be separately estimated alongside saturated month fixed effects. Table~\ref{tab:twoevent} reports the two-event piecewise specification within LLM arrival. Table~\ref{tab:prewindow} reports eight pre-window ASI interactions (joint $F(8)=1.14$, $p=0.33$). Table~\ref{tab:timeplacebo} reports placebo-in-time tests. Table~\ref{tab:lowcopilot} reports H1--H4 outside the first code-LLM languages ($N=1{,}475{,}279$). Table~\ref{tab:tagslope} adds a pre-period tag answer-volume slope control.

\begin{table}[t]
\centering
\caption{Platform-wide secular decline and identification}
\label{tab:secular_control}
\scriptsize
\setlength{\tabcolsep}{3pt}
\begin{tabular}{p{0.28\textwidth}p{0.20\textwidth}p{0.20\textwidth}p{0.24\textwidth}}
\toprule
Specification & Post$\times$ICI & Post$\times$ECI & Treatment of secular decline \\
\midrule
Baseline OLS & $-0.153$ & $0.228$ & Saturated calendar-month fixed effects \\
Add Post $\times$ platform-wide monthly answer volume & Not separately identified & Not separately identified & Collinear with calendar-month fixed effects \\
\bottomrule
\end{tabular}
\end{table}

\begin{table}[t]
\centering
\caption{Two-event piecewise specification within LLM arrival}
\label{tab:twoevent}
\scriptsize
\begin{tabular}{p{0.34\textwidth}p{0.28\textwidth}p{0.28\textwidth}}
\toprule
Window & Joint $F$-test ($p$) & Interpretation \\
\midrule
Code-LLM leads ($k^C=-6$ to $-1$) & $F(6)=0.71$ ($p=0.64$) & Joint test does not reject \\
Code-LLM lags ($k^C=0$ to $+5$) & $F(6)=1.27$ ($p=0.27$) & Short window; test does not reject \\
Leads before public LLM arrival ($k^G=-6$ to $-2$) & Same five-lead test as the event study & Not a second test \\
After public LLM arrival ($k^G=0$ to $+12$) & $F(13)=4.38$ ($p<0.001$) & Jointly significant association \\
\bottomrule
\end{tabular}
\end{table}

\begin{table}[t]
\centering
\caption{Full pre-window $\mathrm{ASI}$ interactions
(net vote score)}
\label{tab:prewindow}
\scriptsize
\begin{tabular}{lcc}
\toprule
Bin, relative to November 2022 & $\beta$ & SE \\
\midrule
January--June 2019 & 0.014 & 0.046 \\
July--December 2019 & $-0.021$ & 0.041 \\
January--June 2020 & 0.008 & 0.037 \\
July--December 2020 & $-0.033$ & 0.034 \\
January--June 2021 & 0.011 & 0.032 \\
July--December 2021 & $-0.027$ & 0.031 \\
January--May 2022 & $-0.019$ & 0.033 \\
June--October 2022 & $-0.046$ & 0.035 \\
\midrule
Joint test of the eight bins & \multicolumn{2}{c}{$F(8)=1.14$, $p=0.33$} \\
\bottomrule
\multicolumn{3}{@{}p{\textwidth}@{}}{\scriptsize The June--October 2022 bin
overlaps early code-LLM diffusion and is the least quiet pre-arrival bin.}
\end{tabular}
\end{table}

\begin{table}[t]
\centering
\caption{Placebo-in-time tests, pre-arrival sample only}
\label{tab:timeplacebo}
\scriptsize
\begin{tabular}{lcccc}
\toprule
Pseudo-event & Post$\times$ICI & SE & Post$\times$ECI & SE \\
\midrule
30 November 2019 & $-0.016$ & 0.047 & 0.024 & 0.061 \\
30 November 2020 & $-0.028$ & 0.043 & 0.019 & 0.056 \\
30 June 2021 & $-0.034$ & 0.041 & 0.038 & 0.053 \\
\midrule
LLM arrival (30 November 2022) & $-0.153$ & 0.036 & 0.228 & 0.041 \\
\bottomrule
\multicolumn{5}{@{}p{\textwidth}@{}}{\scriptsize The last row is the
reported arrival estimate. Each placebo coefficient is smaller than twice
its printed standard error. The largest absolute ratio of coefficient to
standard error is $0.83$. No stored $p$-value is reported for these
cells. The printed standard errors do not indicate a departure at $0.05$,
and the comparison does not prove that the interactions are unique to
late 2022.}
\end{tabular}
\end{table}

\begin{table}[t]
\centering
\caption{H1--H4 estimates outside the first code-LLM languages}
\label{tab:lowcopilot}
\scriptsize
\begin{tabular}{lccc}
\toprule
Coefficient & Reported, full sample & Outside first code-LLM languages & SE \\
\midrule
Post$\times$ICI & $-0.153$ & $-0.086$ & 0.049 \\
Post$\times$ECI & $0.228$ & $0.191$ & 0.064 \\
Post$\times$ASI$\times$ICI & $-0.139$ & $-0.117$ & 0.055 \\
Post$\times$ASI$\times$ECI & $0.176$ & $0.154$ & 0.058 \\
\bottomrule
\multicolumn{4}{@{}p{\textwidth}@{}}{\scriptsize Subsample
$N=1{,}475{,}279$.}
\end{tabular}
\end{table}

\begin{table}[t]
\centering
\caption{Control for pre-period tag answer-volume slopes}
\label{tab:tagslope}
\scriptsize
\begin{tabular}{lcc}
\toprule
Coefficient & Reported & With tag slope \\
\midrule
Post$\times$ICI & $-0.153$ (0.036) & $-0.087$ (0.039) \\
Post$\times$ECI & $0.228$ (0.041) & $0.194$ (0.053) \\
Post$\times$ASI$\times$ICI & $-0.139$ (0.045) & $-0.121$ (0.044) \\
Post$\times$ASI$\times$ECI & $0.176$ (0.052) & $0.149$ (0.046) \\
\bottomrule
\multicolumn{3}{@{}p{\textwidth}@{}}{\scriptsize Parentheses are standard
errors.}
\end{tabular}
\end{table}

\section{Inference, diagnostics, and multiplicity}
\label{app:inference}

Table~\ref{tab:pvalues} reports four-test and eight-test Bonferroni values for the primary net-vote family. Table~\ref{tab:family12} reports the twelve-test Bonferroni family. Table~\ref{tab:diagnostics} summarizes primary model diagnostics. Table~\ref{tab:vif} reports generalized VIFs.

\begin{table}[t]
\centering
\caption{Primary net-vote $p$-values with retrospective multiplicity corrections}
\label{tab:pvalues}
\scriptsize
\begin{tabular}{lccc}
\toprule
H & Raw $p$ & Bonferroni, four tests & Bonferroni, eight-test sensitivity \\
\midrule
H1 & 0.00002 & 0.00008 & 0.00016 \\
H2 & $<$0.0001 & $<$0.0004 & $<$0.0008 \\
H3 & 0.0020 & 0.0080 & 0.016 \\
H4 & 0.00071 & 0.0029 & 0.0057 \\
\bottomrule
\end{tabular}
\end{table}

\begin{table}[t]
\centering
\caption{Twelve-test Bonferroni values computed from printed estimates}
\label{tab:family12}
\scriptsize
\begin{tabular}{llcc}
\toprule
Coefficient & Source $p$ & $\times 12$ & Below 0.05 \\
\midrule
Post$\times$ICI, net votes & 0.00002 & 0.00024 & Yes \\
Post$\times$ECI, net votes & $<$0.0001 & $<$0.0012 & Yes \\
Post$\times$ASI$\times$ICI, net votes & 0.0020 & 0.024 & Yes \\
Post$\times$ASI$\times$ECI, net votes & 0.00071 & 0.0086 & Yes \\
Post$\times$ICI, acceptance & 0.031 & 0.375 & No \\
Post$\times$ECI, acceptance & 0.0003 & 0.004 & Yes \\
Post$\times$ASI$\times$ICI, acceptance & 0.006 & 0.075 & No \\
Post$\times$ASI$\times$ECI, acceptance & 0.0006 & 0.008 & Yes \\
H5, ECI share & $<$0.0001 & $<$0.0012 & Yes \\
H7, selected ASI & $<$0.0001 & $<$0.0012 & Yes \\
H6, volatility & 0.034 & 0.408 & No \\
Post$\times$status, count & 0.0014 & 0.017 & Yes \\
\bottomrule
\multicolumn{4}{@{}p{\textwidth}@{}}{\scriptsize Acceptance and count
$p$-values are two-sided Wald tests computed from the reported
coefficients and standard errors ($z=\hat\beta/\mathrm{SE}$).
H1, H3, and H4 net-vote $p$-values are two-sided normal probabilities
for the printed coefficient divided by its standard error.
H1 uses $t=-0.153/0.036=-4.25$, reported as $0.00002$; twelve times
that reported value is $0.00024$. H3 uses $t=-0.139/0.045=-3.09$
($p=0.0020$). H4 uses $t=0.176/0.052=3.38$ ($p=0.00071$).}
\end{tabular}
\end{table}

\begin{table*}[t]
\centering
\caption{Primary model diagnostics}
\label{tab:diagnostics}
\scriptsize
\setlength{\tabcolsep}{3pt}
\begin{tabular}{p{0.22\textwidth}p{0.17\textwidth}p{0.17\textwidth}p{0.17\textwidth}p{0.17\textwidth}}
\toprule
Diagnostic & Net vote score (OLS) & Acceptance (logit) & ECI share (frac.\ resp.) & Exit hazard (cloglog) \\
\midrule
Estimator & OLS & Logit & Fractional response & Complementary log--log \\
Clustering level & Tag (14,502) & Contributor (354,209) & Contributor & Contributor \\
Outcome mean / variance & 3.14 / 70.9 & 0.42 / 0.24 & 0.38 / 0.06 & Event rate 3.2\% \\
Distribution / zero share & Median 1; zero share 38.2\%; min $-12$, max 1,405 & 0 to 1 & 0 to 1 & 0 to 1 \\
Robust variance estimator & Tag-clustered HC & Contributor-clustered & Contributor-clustered & Contributor-clustered \\
Heteroskedasticity & White test $p<0.001$; tag-clustered robust SE & --- & --- & --- \\
Convergence & Closed-form OLS solution & No perfect prediction; Firth check stable & Converged & Converged \\
Link/specification test & RESET-type $p=0.26$ & Hosmer--Lemeshow $p=0.41$ & RESET-type $p=0.31$ & Time-varying coeff.\ $p=0.34$ \\
Influence / cluster sensitivity & Leave-one-tag-out maximum coefficient change $=0.019$ & --- & --- & --- \\
Calibration (Brier; slope) & --- & 0.21; 0.97 & --- & --- \\
Discrimination & --- & AUC = 0.71 & --- & C-index = 0.74 (attrition) \\
Omnibus Wald test & $p<0.001$ & $p<0.001$ & $p<0.001$ & $p<0.001$ \\
\bottomrule
\end{tabular}
\end{table*}

\begin{table}[t]
\centering
\caption{Generalized VIFs for the valuation specification and the separate supply specification (centered)}
\label{tab:vif}
\scriptsize
\begin{tabular}{lccc}
\toprule
Term & Df & GVIF & GVIF$^{1/(2\,\mathrm{Df})}$ \\
\midrule
\multicolumn{4}{l}{\textit{Valuation specification}} \\
$\mathrm{Post}_t$ & 1 & 1.02 & 1.01 \\
$\mathrm{ICI}_a$ & 1 & 1.18 & 1.09 \\
$\mathrm{ECI}_a$ & 1 & 1.24 & 1.11 \\
$\mathrm{ASI}_{q(a)}$ & 1 & 1.15 & 1.07 \\
$\mathrm{Post}_t\times\mathrm{ICI}_a$ & 1 & 2.31 & 1.52 \\
$\mathrm{Post}_t\times\mathrm{ECI}_a$ & 1 & 2.48 & 1.57 \\
$\mathrm{Post}_t\times\mathrm{ASI}_{q(a)}$ & 1 & 2.12 & 1.46 \\
$\mathrm{Post}_t\times\mathrm{ASI}_{q(a)}\times\mathrm{ICI}_a$ & 1 & 3.85 & 1.96 \\
$\mathrm{Post}_t\times\mathrm{ASI}_{q(a)}\times\mathrm{ECI}_a$ & 1 & 3.92 & 1.98 \\
\midrule
\multicolumn{4}{l}{\textit{Supply specification}} \\
ICI$\times$Post$\times$Supply & 1 & 3.84 & 1.96 \\
ECI$\times$Post$\times$Supply & 1 & 5.29 & 2.30 \\
\bottomrule
\end{tabular}
\par\smallskip
{\scriptsize Note: Df is the number of model-matrix columns assigned to
the term. Each displayed continuous interaction contributes one
column, so the adjusted value is the square root of the reported GVIF.
The two panels are separate specifications.}
\end{table}

\section{Measurement and ASI validation}
\label{app:measurement_ext}

Table~\ref{tab:modelasi} reports agreement between expert ASI and a frozen pre-30~November~2022 model checkpoint on the 600-question expert panel. Correctness rises from $0.19$ in the lowest ASI decile to $0.67$ in the highest; ASI correlates $0.58$ with correctness and $0.51$ with usability. Estimating H3 and H4 on the $N=1{,}380$ attached answers yields $\mathrm{Post}\times\mathrm{Correct}\times\mathrm{ICI}=-0.11$ (SE $0.09$) and $\mathrm{Post}\times\mathrm{Correct}\times\mathrm{ECI}=0.15$ (SE $0.10$), same sign as $-0.139$ and $0.176$ but too imprecise to confirm the answerability interpretation alone. Raising the AI-flag threshold to precision $0.85$ yields Post$\times$ICI $=-0.100$ (SE $0.037$) and Post$\times$ECI $=0.224$ (SE $0.050$), with $N=4{,}153{,}400$.

\begin{table}[t]
\centering
\caption{Agreement between expert ASI and a frozen model
($n=600$ questions)}
\label{tab:modelasi}
\scriptsize
\begin{tabular}{lc}
\toprule
Comparison & Estimate \\
\midrule
ASI with model correctness ($r$) & 0.58 \\
ASI with model usability ($r$) & 0.51 \\
ASI rank with correctness (Spearman) & 0.54 \\
Correctness, ASI decile 1 / decile 10 & 0.19 / 0.67 \\
Usability, ASI decile 1 / decile 10 & 0.24 / 0.61 \\
\bottomrule
\multicolumn{2}{@{}p{\textwidth}@{}}{\scriptsize The 600-question sample is
the expert panel used for ASI validation.}
\end{tabular}
\end{table}

\section{Additional results and sensitivity analyses}
\label{app:extra_results}

Table~\ref{tab:outcomerobust} reports top-1\% trim and ordered outcome specifications. Table~\ref{tab:attrition} and Table~\ref{tab:ipw} document attrition and IPW adaptation estimates. Table~\ref{tab:placebo} reports a placebo outcome. Table~\ref{tab:supply} reports an exploratory informational-supply association (not causal mediation). Table~\ref{tab:exitbounds} reports person-weighted exit scenarios. Table~\ref{tab:deleted} reports hypothetical deleted-answer bounds; these are not estimated regression coefficients and cannot recover deleted posts.

\begin{table*}[t]
\centering
\caption{Robustness of the primary interactions to alternative outcome specifications}
\label{tab:outcomerobust}
\scriptsize
\setlength{\tabcolsep}{4pt}
\begin{tabular}{llccc}
\toprule
Hypothesis & Variable & Baseline (OLS) & Trim top 1\% (OLS) & Ordered outcome (logit) \\
\midrule
H1 & Post$\times$ICI & $-0.153$ (0.036)$^{***}$ & $-0.141$ (0.038)$^{***}$ & $-0.412$ (0.098)$^{***}$ \\
H2 & Post$\times$ECI & $0.228$ (0.041)$^{***}$ & $0.215$ (0.043)$^{***}$ & $0.615$ (0.110)$^{***}$ \\
H3 & Post$\times$ICI$\times$ASI & $-0.139$ (0.045)$^{**}$ & $-0.125$ (0.048)$^{**}$ & $-0.380$ (0.125)$^{**}$ \\
H4 & Post$\times$ECI$\times$ASI & $0.176$ (0.052)$^{***}$ & $0.162$ (0.055)$^{**}$ & $0.475$ (0.140)$^{***}$ \\
\midrule
& Observations & 4,215,082 & 4,172,931 & 4,215,082 \\
& Adjusted $R^2$ / pseudo $R^2$ & 0.142 & 0.138 & 0.095 \\
\bottomrule
\multicolumn{5}{@{}p{\textwidth}@{}}{\scriptsize Standard errors in
parentheses are clustered at the tag level. $^{***}p<0.001$,
$^{**}p<0.01$, from the normal tail of the printed coefficient divided
by its standard error. Baseline cells repeat the net-vote column of
Table~\ref{tab:epistemic}. The trim drops the top 1\% of net-vote
answers. The ordered outcome uses a five-category vote binning. Ordered
coefficients are on the latent index.}
\end{tabular}
\end{table*}

\begin{table}[t]
\centering
\caption{Pre-event characteristics of retained versus exiting contributors}
\label{tab:attrition}
\scriptsize
\begin{tabular}{lcc}
\toprule
Pre-event characteristic & Retained ($n=218{,}447$) & Exited ($n=135{,}762$) \\
\midrule
Median pre-event reputation & 410 & 185 \\
Mean monthly answers (pre) & 2.8 & 1.1 \\
Mean ECI share & 0.38 & 0.31 \\
Mean selected ASI & 0.55 & 0.61 \\
Mean tag entropy & 1.9 & 1.2 \\
\bottomrule
\end{tabular}
\end{table}

\begin{table}[t]
\centering
\caption{Attrition-adjusted adaptation estimates}
\label{tab:ipw}
\scriptsize
\begin{tabular}{lccc}
\toprule
Outcome & Unweighted & IPW & Cohort (explicit zeros) \\
\midrule
ECI share & $0.058^{***}$ & $0.051^{***}$ & $0.047^{**}$ \\
& (0.012) & (0.014) & (0.015) \\
Selected ASI & $-0.082^{***}$ & $-0.074^{***}$ & $-0.069^{**}$ \\
& (0.019) & (0.022) & (0.024) \\
Complexity & $0.045^{**}$ & $0.039^{*}$ & $0.036^{*}$ \\
& (0.015) & (0.017) & (0.018) \\
\bottomrule
\multicolumn{4}{@{}p{\textwidth}@{}}{\scriptsize $^{*}p<0.05$, $^{**}p<0.01$, $^{***}p<0.001$.
The cohort epistemic-share estimate is $0.047/0.015=3.13$ (two-sided $p=0.0017$), so it takes two stars.
The cohort specification retains all 354,209 pre-period contributors.}
\end{tabular}
\end{table}

\begin{table}[t]
\centering
\caption{Placebo test: pre-period edit volume}
\label{tab:placebo}
\scriptsize
\begin{tabular}{lccc}
\toprule
Outcome & Interaction & Estimate & $p$-value \\
\midrule
Pre-period answer edit volume & Post$\times$ASI & 0.004 & 0.612 \\
\bottomrule
\multicolumn{4}{@{}p{\textwidth}@{}}{\scriptsize Note: The coefficient is
interpreted as a placebo differential association, not as proof of a
structural zero.}
\end{tabular}
\end{table}

\begin{table}[t]
\centering
\caption{Informational-supply association}
\label{tab:supply}
\scriptsize
\begin{tabular}{lc}
\toprule
Quantity & Estimate \\
\midrule
Post-period change in supply, high-ASI tags (log points) & 0.17 \\
Post-period change in supply, low-ASI tags (log points) & 0.04 \\
Supply$\times$ECI & 0.092 (SE 0.040) \\
Post$\times$ICI after supply control & $-0.077$ (SE 0.038) \\
Post$\times$ECI after supply control & 0.171 (SE 0.052) \\
Post$\times$ASI$\times$ECI after supply control & 0.128 (SE 0.046) \\
Informational interquartile contrast after supply control & $-0.026$ \\
Share of the $-0.153$ coefficient remaining & 0.503 \\
Share of the $0.228$ coefficient remaining & 0.75 \\
\bottomrule
\multicolumn{2}{@{}p{\textwidth}@{}}{\scriptsize High- and low-ASI tags are
the top and bottom quartiles of mean pre-period ASI. The supply term is
not interpreted as a causal mediator.}
\end{tabular}
\end{table}

\begin{table}[t]
\centering
\caption{Person-weighted scenarios for the status interactions}
\label{tab:exitbounds}
\scriptsize
\begin{tabular}{lcc}
\toprule
Assumption for the 135,762 exiters & H5, ECI share & H7, selected ASI \\
\midrule
Same shift as stayers (reported $0.058$, $-0.082$) & 0.058 & $-0.082$ \\
Half the stayer shift & 0.047 & $-0.066$ \\
No shift & 0.036 & $-0.051$ \\
Opposite shift of the same size & 0.014 & $-0.019$ \\
\bottomrule
\multicolumn{3}{@{}p{\textwidth}@{}}{\scriptsize Weights are $0.617$ on
stayers and $0.383$ on exiters. Calculated from the reported
coefficients and the attrition counts. Not an estimate of exiter
behavior.}
\end{tabular}
\end{table}

\begin{table*}[t]
\centering
\caption{Hypothetical sensitivity bounds under assumed deleted-answer scenarios}
\label{tab:deleted}
\scriptsize
\setlength{\tabcolsep}{3pt}
\begin{tabular}{lcccccc}
\toprule
Scenario & \makecell{Post-period\\deletion} & \makecell{Pre-period\\deletion} & \makecell{Vote penalty\\high ICI} & \makecell{Vote penalty\\high ECI} & \makecell{Hypothetical\\Post$\times$ICI bound} & \makecell{Hypothetical\\Post$\times$ECI bound} \\
\midrule
1 Baseline & 8\% & 3\% & $-1.8$ & $-0.4$ & $-0.119$ & $0.211$ \\
2 Conservative & 5\% & 2\% & $-1.8$ & $-0.4$ & $-0.131$ & $0.218$ \\
3 Aggressive & 10\% & 5\% & $-1.8$ & $-0.4$ & $-0.105$ & $0.202$ \\
4 High penalty & 8\% & 3\% & $-2.5$ & $-0.6$ & $-0.108$ & $0.195$ \\
5 Worst case & 10\% & 5\% & $-2.5$ & $-0.6$ & $-0.092$ & $0.184$ \\
6 Minimal impact & 2\% & 1\% & $-1.0$ & $-0.2$ & $-0.142$ & $0.225$ \\
\bottomrule
\multicolumn{7}{@{}p{\textwidth}@{}}{\scriptsize These values are
hypothetical bounds from applying assumed vote penalties to assumed
deletion rates. They are not estimated regression coefficients, and they
cannot be reproduced from the exported covariance matrices. The signs
match the reported $-0.153$ and $0.228$ in every row. The bounds
illustrate directional robustness under those assumptions. They do not
recover the deleted posts.}
\end{tabular}
\end{table*}
.

\section{Sample construction details}
\label{app:sample}

Table~\ref{tab:sampleledger} records verified analysis-sample counts. Table~\ref{tab:sensitivitysamples} reports H1--H4 coefficients under alternative sample restrictions. Under those alternative exclusions (community-wiki, self-answers, AI-flagged, and all three), the informational coefficient lies between $-0.096$ and $-0.104$ (maximum absolute shift $\Delta\beta_3=0.057$ from baseline $-0.153$). The epistemic shift remains at most $\Delta\beta_4=0.013$. The ``Exclude AI-flagged'' row uses a noisy text-classification proxy (precision 0.74; recall 0.81), not verified authorship.

\textbf{Automated-account rule.} An account was classified as likely automated if it satisfied at least two of three criteria: (i)~more than 50 answers posted within any 24-hour window on three or more separate days; (ii)~pairwise Jaccard similarity $>0.95$ across $\geq$80\% of its answer texts; or (iii)~account creation and first answer within fewer than 90 seconds with no subsequent editing activity. The rule removed 47,218 answers attributed to 3,841 accounts. Manual inspection of 200 randomly sampled flagged accounts confirmed that 194 exhibited bot-like patterns; the remaining 6 were high-volume human contributors. Sensitivity analyses retaining these 6 accounts do not alter the signs or inference of the primary signed-score OLS interactions.

\textbf{Duplicate handling.} Questions linked as duplicates by community closure were retained with a binary duplicate-link indicator included in $\mathbf{X}_{at}$. Exact-text near-duplicates (Jaccard $>0.98$) posted to the \emph{same} question were removed (12,847 records); near-duplicates across different questions were retained.

\begin{table}[t]
\centering
\caption{Verified analysis sample ledger}
\label{tab:sampleledger}
\scriptsize
\begin{tabular}{lr}
\toprule
Field & Verified value \\
\midrule
\texttt{final\_answer\_rows} & 4,215,082 \\
\texttt{unique\_questions} & 1,842,110 \\
\texttt{unique\_contributors} & 354,209 \\
\texttt{post\_answers\_jan23\_dec25} & 1,712,539 \\
\texttt{pre\_and\_transition\_answers} & 2,502,543 \\
\texttt{ai\_flagged\_post\_answers} & 94,217 \\
\texttt{post\_flag\_prevalence} & 5.50\% \\
\texttt{sensitivity\_sample\_exclude\_ai\_flagged} & 4,120,865 \\
\bottomrule
\end{tabular}
\end{table}

\begin{table}[t]
\centering
\caption{Sensitivity of H1--H4 coefficients to sample restrictions}
\label{tab:sensitivitysamples}
\scriptsize
\begin{tabular}{lccc}
\toprule
Restriction & $\beta_3$ (Post$\times$ICI) & $\beta_4$ (Post$\times$ECI) & $N$ \\
\midrule
Baseline (all retained) & $-0.153^{***}$ & $0.228^{***}$ & 4,215,082 \\
Exclude community-wiki & $-0.100^{**}$ & $0.224^{***}$ & 4,127,670 \\
Exclude self-answers & $-0.098^{**}$ & $0.222^{***}$ & 3,791,226 \\
Exclude AI-flagged & $-0.104^{**}$ & $0.231^{***}$ & 4,120,865 \\
Exclude all three & $-0.096^{*}$ & $0.218^{***}$ & 3,703,449 \\
\bottomrule
\multicolumn{4}{@{}p{\textwidth}@{}}{\scriptsize Note: All models include full
controls and fixed effects. ``Exclude AI-flagged'' removes answers selected
by a noisy text-classification proxy (precision 0.74; recall 0.81), not a
verified set of AI-authored answers. The row is therefore a perturbation
test, not evidence that the result holds in an AI-free sample.}
\end{tabular}
\end{table}

\section{AI-assisted text detector}
\label{app:detector}

Likely AI-assisted text was identified using a fine-tuned DeBERTa-v3-large classifier (checkpoint \texttt{microsoft/deberta-v3-large}, Hugging Face transformers 4.38.1) with a binary classification head. The model was trained on 18,000 labeled examples: 9,000 human-written answers sampled from January~2019--October~2022 and 9,000 synthetic answers generated by prompting GPT-3.5-Turbo (March~2023 snapshot) and GPT-4 (August~2023 snapshot) with 4,500 held-out Stack Overflow questions not present in the analytical sample. The positive class is defined as probability $\geq 0.72$, selected by maximizing Youden's $J$ on a 3,000-item validation partition. Training human texts are drawn exclusively from the pre-event window; classification targets are post-January~2023 answers.

Three trained annotators, blinded to detector scores, rated a stratified sample of 1,000 answers. The majority-vote reference is positive for 280 answers and negative for 720. Table~\ref{tab:detectorperf} reports performance. Cell counts are 227 true positives, 53 false negatives, 79 false positives, and 641 true negatives. Performance across sub-periods (2023--2025) showed no statistically significant decline in AUC (range 0.89--0.92; likelihood-ratio test $\chi^2(2)=3.1$, $p=0.21$). Of 1,712,539 post-period answers, 94,217 (5.50\%) exceeded the 0.72 threshold. This figure is an apparent classification rate, not an estimate of true AI-text prevalence. The flag defines sensitivity samples only and is not interpreted as verified authorship.

\begin{table}[t]
\centering
\caption{AI-assisted text detector performance ($n=1{,}000$ human-validation sample)}
\label{tab:detectorperf}
\scriptsize
\begin{tabular}{lc}
\toprule
Metric & Value (95\% CI) \\
\midrule
Sensitivity (recall) & 0.81 [0.76, 0.86] \\
Specificity & 0.89 [0.87, 0.91] \\
Precision & 0.74 [0.69, 0.79] \\
F1-score & 0.77 \\
AUC-ROC & 0.91 [0.89, 0.93] \\
Reference-positive rate & 28\% (280 / 1,000) \\
\bottomrule
\multicolumn{2}{@{}p{\textwidth}@{}}{\scriptsize Counts are 227, 53, 79,
and 641. Sensitivity is $227/280=0.81$, specificity is $641/720=0.89$,
and precision is $227/(227+79)=0.74$. The same precision is
$(0.81\times 0.28)/[(0.81\times 0.28)+(0.11\times 0.72)]=0.74$.
F1 is $2\times 227/(2\times 227+79+53)=0.77$. Wald intervals use those
cell counts.}
\end{tabular}
\end{table}

\section{Notation and coding anchors}
\label{app:notation}

Table~\ref{tab:notation} records principal indices and model notation. Table~\ref{tab:codinganchors} reports coding anchors for informational and epistemic contribution.

\begin{table}[t]
\centering
\caption{Principal indices and model notation}
\label{tab:notation}
\scriptsize
\begin{tabular}{@{}p{0.16\textwidth}p{0.46\textwidth}p{0.26\textwidth}@{}}
\toprule
Symbol & Standardized meaning & Primary use \\
\midrule
$a$ & Answer; the focal answer-level observation & H1--H4 \\
$q$; $q(a)$ & Question containing answer $a$ & ASI, complexity \\
$u$; $u(a)$ & Contributor; author of answer $a$ & H5/H7 panel \\
$t$ & Calendar month & Post indicator, FEs \\
$k$ & Event month relative to focal event & Event study \\
$g$ & Contribution function $\in\{\mathrm{I},\mathrm{E},\mathrm{C}\}$ & Theory only \\
$s(a)$ & Fixed-effect stratum: primary-tag $\times$ question-creation year & Eq.~\eqref{eq:did} \\
$Y_{at}$; $Y_{ut}$ & Answer-level; contributor-month outcome & Outcomes \\
$\mathrm{ICI}_a$; $\mathrm{ECI}_a$ & Informational; epistemic index & Valuation \\
$\mathrm{ASI}_q$ & Expert-rated AI answerability of question $q$ & Exposure proxy \\
$\mathrm{Post}_t$ & Post-event month indicator & Change specs \\
$\mathrm{Status}^{\mathrm{pre}}_u$ & Standardized log pre-event reputation & Adaptation \\
$\mathbf{X}_{at}$; $\mathbf{Z}_{ut}$ & Answer-level; contributor-month controls & Covariates \\
$\mathrm{Volatility}_{\operatorname{tag}(q),t}$ & Tag-level volatility & H6 \\
$\mathrm{RV}_{g,t}(\mathrm{ASI}_q)$ & Return to function $g$ & Theory \\
$\beta$; $\beta_k$; $\beta_1$--$\beta_4$; $\rho$ & Parameters & Estimands \\
$\theta_{\mathrm{I}}$; $\theta_{\mathrm{E}}$ & Triple-interaction coefficients & H3, H4 \\
$\boldsymbol{\gamma}$; $\boldsymbol{\kappa}$ & Control coefficient vectors & Controls \\
$\alpha_{s(a)}$; $\delta_t$ & Stratum; month fixed effects & Baseline \\
$\mu_{u(a)}$; $\lambda_{\operatorname{tag}(q(a))}$; $\eta_u$ & Contributor; tag; panel FEs & Valuation/adaptation \\
$\varepsilon_{at}$; $\varepsilon_{ut}$ & Disturbances & Errors \\
$\mathrm{EventTime}_t$; $\mathbb{1}\{\cdot\}$ & Event month; indicator & Event study \\
\bottomrule
\end{tabular}
\end{table}

\begin{table*}[t]
\centering
\caption{Coding anchors for informational and epistemic contribution}
\label{tab:codinganchors}
\scriptsize
\begin{tabular}{p{0.18\textwidth}p{0.34\textwidth}p{0.34\textwidth}}
\toprule
Pattern & Coding anchor & Illustrative answer form \\
\midrule
High ICI, low ECI & Supplies a direct command, code fragment, or standard
procedure without visible tests, assumptions, or limits & ``Use function
X with argument Y.'' \\
Low ICI, high ECI & Diagnoses uncertainty, compares explanations, or
specifies tests and boundary conditions without offering a complete
drop-in solution & ``This depends on the runtime version; test A before
choosing between B and C.'' \\
High ICI, high ECI & Supplies a usable solution and makes its warrant,
scope, and failure conditions inspectable & ``Use X; the example verifies
condition Y, but it fails under Z and requires version V.'' \\
Low ICI, low ECI & Provides neither a usable solution nor inspectable
reasoning & Unsupported assertion, social comment, or restatement of the
question \\
\bottomrule
\end{tabular}
\end{table*}

\section{Identification and robustness tables}
\label{app:identification}

Table~\ref{tab:secular_control} records that a $\mathrm{Post}_t\times\mathrm{PlatformAnswerVolume}_t$ term cannot be separately estimated alongside saturated month fixed effects. Table~\ref{tab:twoevent} reports the two-event piecewise specification within LLM arrival. Table~\ref{tab:prewindow} reports eight pre-window ASI interactions (joint $F(8)=1.14$, $p=0.33$). Table~\ref{tab:timeplacebo} reports placebo-in-time tests. Table~\ref{tab:lowcopilot} reports H1--H4 outside the first code-LLM languages ($N=1{,}475{,}279$). Table~\ref{tab:tagslope} adds a pre-period tag answer-volume slope control.

\begin{table}[t]
\centering
\caption{Platform-wide secular decline and identification}
\label{tab:secular_control}
\scriptsize
\setlength{\tabcolsep}{3pt}
\begin{tabular}{p{0.28\textwidth}p{0.20\textwidth}p{0.20\textwidth}p{0.24\textwidth}}
\toprule
Specification & Post$\times$ICI & Post$\times$ECI & Treatment of secular decline \\
\midrule
Baseline OLS & $-0.153$ & $0.228$ & Saturated calendar-month fixed effects \\
Add Post $\times$ platform-wide monthly answer volume & Not separately identified & Not separately identified & Collinear with calendar-month fixed effects \\
\bottomrule
\end{tabular}
\end{table}

\begin{table}[t]
\centering
\caption{Two-event piecewise specification within LLM arrival}
\label{tab:twoevent}
\scriptsize
\begin{tabular}{p{0.34\textwidth}p{0.28\textwidth}p{0.28\textwidth}}
\toprule
Window & Joint $F$-test ($p$) & Interpretation \\
\midrule
Code-LLM leads ($k^C=-6$ to $-1$) & $F(6)=0.71$ ($p=0.64$) & Joint test does not reject \\
Code-LLM lags ($k^C=0$ to $+5$) & $F(6)=1.27$ ($p=0.27$) & Short window; test does not reject \\
Leads before public LLM arrival ($k^G=-6$ to $-2$) & Same five-lead test as the event study & Not a second test \\
After public LLM arrival ($k^G=0$ to $+12$) & $F(13)=4.38$ ($p<0.001$) & Jointly significant association \\
\bottomrule
\end{tabular}
\end{table}

\begin{table}[t]
\centering
\caption{Full pre-window $\mathrm{ASI}$ interactions
(net vote score)}
\label{tab:prewindow}
\scriptsize
\begin{tabular}{lcc}
\toprule
Bin, relative to November 2022 & $\beta$ & SE \\
\midrule
January--June 2019 & 0.014 & 0.046 \\
July--December 2019 & $-0.021$ & 0.041 \\
January--June 2020 & 0.008 & 0.037 \\
July--December 2020 & $-0.033$ & 0.034 \\
January--June 2021 & 0.011 & 0.032 \\
July--December 2021 & $-0.027$ & 0.031 \\
January--May 2022 & $-0.019$ & 0.033 \\
June--October 2022 & $-0.046$ & 0.035 \\
\midrule
Joint test of the eight bins & \multicolumn{2}{c}{$F(8)=1.14$, $p=0.33$} \\
\bottomrule
\multicolumn{3}{@{}p{\textwidth}@{}}{\scriptsize The June--October 2022 bin
overlaps early code-LLM diffusion and is the least quiet pre-arrival bin.}
\end{tabular}
\end{table}

\begin{table}[t]
\centering
\caption{Placebo-in-time tests, pre-arrival sample only}
\label{tab:timeplacebo}
\scriptsize
\begin{tabular}{lcccc}
\toprule
Pseudo-event & Post$\times$ICI & SE & Post$\times$ECI & SE \\
\midrule
30 November 2019 & $-0.016$ & 0.047 & 0.024 & 0.061 \\
30 November 2020 & $-0.028$ & 0.043 & 0.019 & 0.056 \\
30 June 2021 & $-0.034$ & 0.041 & 0.038 & 0.053 \\
\midrule
LLM arrival (30 November 2022) & $-0.153$ & 0.036 & 0.228 & 0.041 \\
\bottomrule
\multicolumn{5}{@{}p{\textwidth}@{}}{\scriptsize The last row is the
reported arrival estimate. Each placebo coefficient is smaller than twice
its printed standard error. The largest absolute ratio of coefficient to
standard error is $0.83$. No stored $p$-value is reported for these
cells. The printed standard errors do not indicate a departure at $0.05$,
and the comparison does not prove that the interactions are unique to
late 2022.}
\end{tabular}
\end{table}

\begin{table}[t]
\centering
\caption{H1--H4 estimates outside the first code-LLM languages}
\label{tab:lowcopilot}
\scriptsize
\begin{tabular}{lccc}
\toprule
Coefficient & Reported, full sample & Outside first code-LLM languages & SE \\
\midrule
Post$\times$ICI & $-0.153$ & $-0.086$ & 0.049 \\
Post$\times$ECI & $0.228$ & $0.191$ & 0.064 \\
Post$\times$ASI$\times$ICI & $-0.139$ & $-0.117$ & 0.055 \\
Post$\times$ASI$\times$ECI & $0.176$ & $0.154$ & 0.058 \\
\bottomrule
\multicolumn{4}{@{}p{\textwidth}@{}}{\scriptsize Subsample
$N=1{,}475{,}279$.}
\end{tabular}
\end{table}

\begin{table}[t]
\centering
\caption{Control for pre-period tag answer-volume slopes}
\label{tab:tagslope}
\scriptsize
\begin{tabular}{lcc}
\toprule
Coefficient & Reported & With tag slope \\
\midrule
Post$\times$ICI & $-0.153$ (0.036) & $-0.087$ (0.039) \\
Post$\times$ECI & $0.228$ (0.041) & $0.194$ (0.053) \\
Post$\times$ASI$\times$ICI & $-0.139$ (0.045) & $-0.121$ (0.044) \\
Post$\times$ASI$\times$ECI & $0.176$ (0.052) & $0.149$ (0.046) \\
\bottomrule
\multicolumn{3}{@{}p{\textwidth}@{}}{\scriptsize Parentheses are standard
errors.}
\end{tabular}
\end{table}

\section{Inference, diagnostics, and multiplicity}
\label{app:inference}

Table~\ref{tab:pvalues} reports four-test and eight-test Bonferroni values for the primary net-vote family. Table~\ref{tab:family12} reports the twelve-test Bonferroni family. Table~\ref{tab:diagnostics} summarizes primary model diagnostics. Table~\ref{tab:vif} reports generalized VIFs.

\begin{table}[t]
\centering
\caption{Primary net-vote $p$-values with retrospective multiplicity corrections}
\label{tab:pvalues}
\scriptsize
\begin{tabular}{lccc}
\toprule
H & Raw $p$ & Bonferroni, four tests & Bonferroni, eight-test sensitivity \\
\midrule
H1 & 0.00002 & 0.00008 & 0.00016 \\
H2 & $<$0.0001 & $<$0.0004 & $<$0.0008 \\
H3 & 0.0020 & 0.0080 & 0.016 \\
H4 & 0.00071 & 0.0029 & 0.0057 \\
\bottomrule
\end{tabular}
\end{table}

\begin{table}[t]
\centering
\caption{Twelve-test Bonferroni values computed from printed estimates}
\label{tab:family12}
\scriptsize
\begin{tabular}{llcc}
\toprule
Coefficient & Source $p$ & $\times 12$ & Below 0.05 \\
\midrule
Post$\times$ICI, net votes & 0.00002 & 0.00024 & Yes \\
Post$\times$ECI, net votes & $<$0.0001 & $<$0.0012 & Yes \\
Post$\times$ASI$\times$ICI, net votes & 0.0020 & 0.024 & Yes \\
Post$\times$ASI$\times$ECI, net votes & 0.00071 & 0.0086 & Yes \\
Post$\times$ICI, acceptance & 0.031 & 0.375 & No \\
Post$\times$ECI, acceptance & 0.0003 & 0.004 & Yes \\
Post$\times$ASI$\times$ICI, acceptance & 0.006 & 0.075 & No \\
Post$\times$ASI$\times$ECI, acceptance & 0.0006 & 0.008 & Yes \\
H5, ECI share & $<$0.0001 & $<$0.0012 & Yes \\
H7, selected ASI & $<$0.0001 & $<$0.0012 & Yes \\
H6, volatility & 0.034 & 0.408 & No \\
Post$\times$status, count & 0.0014 & 0.017 & Yes \\
\bottomrule
\multicolumn{4}{@{}p{\textwidth}@{}}{\scriptsize Acceptance and count
$p$-values are two-sided Wald tests computed from the reported
coefficients and standard errors ($z=\hat\beta/\mathrm{SE}$).
H1, H3, and H4 net-vote $p$-values are two-sided normal probabilities
for the printed coefficient divided by its standard error.
H1 uses $t=-0.153/0.036=-4.25$, reported as $0.00002$; twelve times
that reported value is $0.00024$. H3 uses $t=-0.139/0.045=-3.09$
($p=0.0020$). H4 uses $t=0.176/0.052=3.38$ ($p=0.00071$).}
\end{tabular}
\end{table}

\begin{table*}[t]
\centering
\caption{Primary model diagnostics}
\label{tab:diagnostics}
\scriptsize
\setlength{\tabcolsep}{3pt}
\begin{tabular}{p{0.22\textwidth}p{0.17\textwidth}p{0.17\textwidth}p{0.17\textwidth}p{0.17\textwidth}}
\toprule
Diagnostic & Net vote score (OLS) & Acceptance (logit) & ECI share (frac.\ resp.) & Exit hazard (cloglog) \\
\midrule
Estimator & OLS & Logit & Fractional response & Complementary log--log \\
Clustering level & Tag (14,502) & Contributor (354,209) & Contributor & Contributor \\
Outcome mean / variance & 3.14 / 70.9 & 0.42 / 0.24 & 0.38 / 0.06 & Event rate 3.2\% \\
Distribution / zero share & Median 1; zero share 38.2\%; min $-12$, max 1,405 & 0 to 1 & 0 to 1 & 0 to 1 \\
Robust variance estimator & Tag-clustered HC & Contributor-clustered & Contributor-clustered & Contributor-clustered \\
Heteroskedasticity & White test $p<0.001$; tag-clustered robust SE & --- & --- & --- \\
Convergence & Closed-form OLS solution & No perfect prediction; Firth check stable & Converged & Converged \\
Link/specification test & RESET-type $p=0.26$ & Hosmer--Lemeshow $p=0.41$ & RESET-type $p=0.31$ & Time-varying coeff.\ $p=0.34$ \\
Influence / cluster sensitivity & Leave-one-tag-out maximum coefficient change $=0.019$ & --- & --- & --- \\
Calibration (Brier; slope) & --- & 0.21; 0.97 & --- & --- \\
Discrimination & --- & AUC = 0.71 & --- & C-index = 0.74 (attrition) \\
Omnibus Wald test & $p<0.001$ & $p<0.001$ & $p<0.001$ & $p<0.001$ \\
\bottomrule
\end{tabular}
\end{table*}

\begin{table}[t]
\centering
\caption{Generalized VIFs for the valuation specification and the separate supply specification (centered)}
\label{tab:vif}
\scriptsize
\begin{tabular}{lccc}
\toprule
Term & Df & GVIF & GVIF$^{1/(2\,\mathrm{Df})}$ \\
\midrule
\multicolumn{4}{l}{\textit{Valuation specification}} \\
$\mathrm{Post}_t$ & 1 & 1.02 & 1.01 \\
$\mathrm{ICI}_a$ & 1 & 1.18 & 1.09 \\
$\mathrm{ECI}_a$ & 1 & 1.24 & 1.11 \\
$\mathrm{ASI}_{q(a)}$ & 1 & 1.15 & 1.07 \\
$\mathrm{Post}_t\times\mathrm{ICI}_a$ & 1 & 2.31 & 1.52 \\
$\mathrm{Post}_t\times\mathrm{ECI}_a$ & 1 & 2.48 & 1.57 \\
$\mathrm{Post}_t\times\mathrm{ASI}_{q(a)}$ & 1 & 2.12 & 1.46 \\
$\mathrm{Post}_t\times\mathrm{ASI}_{q(a)}\times\mathrm{ICI}_a$ & 1 & 3.85 & 1.96 \\
$\mathrm{Post}_t\times\mathrm{ASI}_{q(a)}\times\mathrm{ECI}_a$ & 1 & 3.92 & 1.98 \\
\midrule
\multicolumn{4}{l}{\textit{Supply specification}} \\
ICI$\times$Post$\times$Supply & 1 & 3.84 & 1.96 \\
ECI$\times$Post$\times$Supply & 1 & 5.29 & 2.30 \\
\bottomrule
\end{tabular}
\par\smallskip
{\scriptsize Note: Df is the number of model-matrix columns assigned to
the term. Each displayed continuous interaction contributes one
column, so the adjusted value is the square root of the reported GVIF.
The two panels are separate specifications.}
\end{table}

\section{Measurement and ASI validation}
\label{app:measurement_ext}

Table~\ref{tab:modelasi} reports agreement between expert ASI and a frozen pre-30~November~2022 model checkpoint on the 600-question expert panel. Correctness rises from $0.19$ in the lowest ASI decile to $0.67$ in the highest; ASI correlates $0.58$ with correctness and $0.51$ with usability. Estimating H3 and H4 on the $N=1{,}380$ attached answers yields $\mathrm{Post}\times\mathrm{Correct}\times\mathrm{ICI}=-0.11$ (SE $0.09$) and $\mathrm{Post}\times\mathrm{Correct}\times\mathrm{ECI}=0.15$ (SE $0.10$), same sign as $-0.139$ and $0.176$ but too imprecise to confirm the answerability interpretation alone. Raising the AI-flag threshold to precision $0.85$ yields Post$\times$ICI $=-0.100$ (SE $0.037$) and Post$\times$ECI $=0.224$ (SE $0.050$), with $N=4{,}153{,}400$.

\begin{table}[t]
\centering
\caption{Agreement between expert ASI and a frozen model
($n=600$ questions)}
\label{tab:modelasi}
\scriptsize
\begin{tabular}{lc}
\toprule
Comparison & Estimate \\
\midrule
ASI with model correctness ($r$) & 0.58 \\
ASI with model usability ($r$) & 0.51 \\
ASI rank with correctness (Spearman) & 0.54 \\
Correctness, ASI decile 1 / decile 10 & 0.19 / 0.67 \\
Usability, ASI decile 1 / decile 10 & 0.24 / 0.61 \\
\bottomrule
\multicolumn{2}{@{}p{\textwidth}@{}}{\scriptsize The 600-question sample is
the expert panel used for ASI validation.}
\end{tabular}
\end{table}

\section{Additional results and sensitivity analyses}
\label{app:extra_results}

Table~\ref{tab:outcomerobust} reports top-1\% trim and ordered outcome specifications. Table~\ref{tab:attrition} and Table~\ref{tab:ipw} document attrition and IPW adaptation estimates. Table~\ref{tab:placebo} reports a placebo outcome. Table~\ref{tab:supply} reports an exploratory informational-supply association (not causal mediation). Table~\ref{tab:exitbounds} reports person-weighted exit scenarios. Table~\ref{tab:deleted} reports hypothetical deleted-answer bounds; these are not estimated regression coefficients and cannot recover deleted posts.

\begin{table*}[t]
\centering
\caption{Robustness of the primary interactions to alternative outcome specifications}
\label{tab:outcomerobust}
\scriptsize
\setlength{\tabcolsep}{4pt}
\begin{tabular}{llccc}
\toprule
Hypothesis & Variable & Baseline (OLS) & Trim top 1\% (OLS) & Ordered outcome (logit) \\
\midrule
H1 & Post$\times$ICI & $-0.153$ (0.036)$^{***}$ & $-0.141$ (0.038)$^{***}$ & $-0.412$ (0.098)$^{***}$ \\
H2 & Post$\times$ECI & $0.228$ (0.041)$^{***}$ & $0.215$ (0.043)$^{***}$ & $0.615$ (0.110)$^{***}$ \\
H3 & Post$\times$ICI$\times$ASI & $-0.139$ (0.045)$^{**}$ & $-0.125$ (0.048)$^{**}$ & $-0.380$ (0.125)$^{**}$ \\
H4 & Post$\times$ECI$\times$ASI & $0.176$ (0.052)$^{***}$ & $0.162$ (0.055)$^{**}$ & $0.475$ (0.140)$^{***}$ \\
\midrule
& Observations & 4,215,082 & 4,172,931 & 4,215,082 \\
& Adjusted $R^2$ / pseudo $R^2$ & 0.142 & 0.138 & 0.095 \\
\bottomrule
\multicolumn{5}{@{}p{\textwidth}@{}}{\scriptsize Standard errors in
parentheses are clustered at the tag level. $^{***}p<0.001$,
$^{**}p<0.01$, from the normal tail of the printed coefficient divided
by its standard error. Baseline cells repeat the net-vote column of
Table~\ref{tab:epistemic}. The trim drops the top 1\% of net-vote
answers. The ordered outcome uses a five-category vote binning. Ordered
coefficients are on the latent index.}
\end{tabular}
\end{table*}

\begin{table}[t]
\centering
\caption{Pre-event characteristics of retained versus exiting contributors}
\label{tab:attrition}
\scriptsize
\begin{tabular}{lcc}
\toprule
Pre-event characteristic & Retained ($n=218{,}447$) & Exited ($n=135{,}762$) \\
\midrule
Median pre-event reputation & 410 & 185 \\
Mean monthly answers (pre) & 2.8 & 1.1 \\
Mean ECI share & 0.38 & 0.31 \\
Mean selected ASI & 0.55 & 0.61 \\
Mean tag entropy & 1.9 & 1.2 \\
\bottomrule
\end{tabular}
\end{table}

\begin{table}[t]
\centering
\caption{Attrition-adjusted adaptation estimates}
\label{tab:ipw}
\scriptsize
\begin{tabular}{lccc}
\toprule
Outcome & Unweighted & IPW & Cohort (explicit zeros) \\
\midrule
ECI share & $0.058^{***}$ & $0.051^{***}$ & $0.047^{**}$ \\
& (0.012) & (0.014) & (0.015) \\
Selected ASI & $-0.082^{***}$ & $-0.074^{***}$ & $-0.069^{**}$ \\
& (0.019) & (0.022) & (0.024) \\
Complexity & $0.045^{**}$ & $0.039^{*}$ & $0.036^{*}$ \\
& (0.015) & (0.017) & (0.018) \\
\bottomrule
\multicolumn{4}{@{}p{\textwidth}@{}}{\scriptsize $^{*}p<0.05$, $^{**}p<0.01$, $^{***}p<0.001$.
The cohort epistemic-share estimate is $0.047/0.015=3.13$ (two-sided $p=0.0017$), so it takes two stars.
The cohort specification retains all 354,209 pre-period contributors.}
\end{tabular}
\end{table}

\begin{table}[t]
\centering
\caption{Placebo test: pre-period edit volume}
\label{tab:placebo}
\scriptsize
\begin{tabular}{lccc}
\toprule
Outcome & Interaction & Estimate & $p$-value \\
\midrule
Pre-period answer edit volume & Post$\times$ASI & 0.004 & 0.612 \\
\bottomrule
\multicolumn{4}{@{}p{\textwidth}@{}}{\scriptsize Note: The coefficient is
interpreted as a placebo differential association, not as proof of a
structural zero.}
\end{tabular}
\end{table}

\begin{table}[t]
\centering
\caption{Informational-supply association}
\label{tab:supply}
\scriptsize
\begin{tabular}{lc}
\toprule
Quantity & Estimate \\
\midrule
Post-period change in supply, high-ASI tags (log points) & 0.17 \\
Post-period change in supply, low-ASI tags (log points) & 0.04 \\
Supply$\times$ECI & 0.092 (SE 0.040) \\
Post$\times$ICI after supply control & $-0.077$ (SE 0.038) \\
Post$\times$ECI after supply control & 0.171 (SE 0.052) \\
Post$\times$ASI$\times$ECI after supply control & 0.128 (SE 0.046) \\
Informational interquartile contrast after supply control & $-0.026$ \\
Share of the $-0.153$ coefficient remaining & 0.503 \\
Share of the $0.228$ coefficient remaining & 0.75 \\
\bottomrule
\multicolumn{2}{@{}p{\textwidth}@{}}{\scriptsize High- and low-ASI tags are
the top and bottom quartiles of mean pre-period ASI. The supply term is
not interpreted as a causal mediator.}
\end{tabular}
\end{table}

\begin{table}[t]
\centering
\caption{Person-weighted scenarios for the status interactions}
\label{tab:exitbounds}
\scriptsize
\begin{tabular}{lcc}
\toprule
Assumption for the 135,762 exiters & H5, ECI share & H7, selected ASI \\
\midrule
Same shift as stayers (reported $0.058$, $-0.082$) & 0.058 & $-0.082$ \\
Half the stayer shift & 0.047 & $-0.066$ \\
No shift & 0.036 & $-0.051$ \\
Opposite shift of the same size & 0.014 & $-0.019$ \\
\bottomrule
\multicolumn{3}{@{}p{\textwidth}@{}}{\scriptsize Weights are $0.617$ on
stayers and $0.383$ on exiters. Calculated from the reported
coefficients and the attrition counts. Not an estimate of exiter
behavior.}
\end{tabular}
\end{table}

\begin{table*}[t]
\centering
\caption{Hypothetical sensitivity bounds under assumed deleted-answer scenarios}
\label{tab:deleted}
\scriptsize
\setlength{\tabcolsep}{3pt}
\begin{tabular}{lcccccc}
\toprule
Scenario & \makecell{Post-period\\deletion} & \makecell{Pre-period\\deletion} & \makecell{Vote penalty\\high ICI} & \makecell{Vote penalty\\high ECI} & \makecell{Hypothetical\\Post$\times$ICI bound} & \makecell{Hypothetical\\Post$\times$ECI bound} \\
\midrule
1 Baseline & 8\% & 3\% & $-1.8$ & $-0.4$ & $-0.119$ & $0.211$ \\
2 Conservative & 5\% & 2\% & $-1.8$ & $-0.4$ & $-0.131$ & $0.218$ \\
3 Aggressive & 10\% & 5\% & $-1.8$ & $-0.4$ & $-0.105$ & $0.202$ \\
4 High penalty & 8\% & 3\% & $-2.5$ & $-0.6$ & $-0.108$ & $0.195$ \\
5 Worst case & 10\% & 5\% & $-2.5$ & $-0.6$ & $-0.092$ & $0.184$ \\
6 Minimal impact & 2\% & 1\% & $-1.0$ & $-0.2$ & $-0.142$ & $0.225$ \\
\bottomrule
\multicolumn{7}{@{}p{\textwidth}@{}}{\scriptsize These values are
hypothetical bounds from applying assumed vote penalties to assumed
deletion rates. They are not estimated regression coefficients, and they
cannot be reproduced from the exported covariance matrices. The signs
match the reported $-0.153$ and $0.228$ in every row. The bounds
illustrate directional robustness under those assumptions. They do not
recover the deleted posts.}
\end{tabular}
\end{table*}
.

\section{Sample construction details}
\label{app:sample}

Table~\ref{tab:sampleledger} records verified analysis-sample counts. Table~\ref{tab:sensitivitysamples} reports H1--H4 coefficients under alternative sample restrictions. Under those alternative exclusions (community-wiki, self-answers, AI-flagged, and all three), the informational coefficient lies between $-0.096$ and $-0.104$ (maximum absolute shift $\Delta\beta_3=0.057$ from baseline $-0.153$). The epistemic shift remains at most $\Delta\beta_4=0.013$. The ``Exclude AI-flagged'' row uses a noisy text-classification proxy (precision 0.74; recall 0.81), not verified authorship.

\textbf{Automated-account rule.} An account was classified as likely automated if it satisfied at least two of three criteria: (i)~more than 50 answers posted within any 24-hour window on three or more separate days; (ii)~pairwise Jaccard similarity $>0.95$ across $\geq$80\% of its answer texts; or (iii)~account creation and first answer within fewer than 90 seconds with no subsequent editing activity. The rule removed 47,218 answers attributed to 3,841 accounts. Manual inspection of 200 randomly sampled flagged accounts confirmed that 194 exhibited bot-like patterns; the remaining 6 were high-volume human contributors. Sensitivity analyses retaining these 6 accounts do not alter the signs or inference of the primary signed-score OLS interactions.

\textbf{Duplicate handling.} Questions linked as duplicates by community closure were retained with a binary duplicate-link indicator included in $\mathbf{X}_{at}$. Exact-text near-duplicates (Jaccard $>0.98$) posted to the \emph{same} question were removed (12,847 records); near-duplicates across different questions were retained.

\begin{table}[t]
\centering
\caption{Verified analysis sample ledger}
\label{tab:sampleledger}
\scriptsize
\begin{tabular}{lr}
\toprule
Field & Verified value \\
\midrule
\texttt{final\_answer\_rows} & 4,215,082 \\
\texttt{unique\_questions} & 1,842,110 \\
\texttt{unique\_contributors} & 354,209 \\
\texttt{post\_answers\_jan23\_dec25} & 1,712,539 \\
\texttt{pre\_and\_transition\_answers} & 2,502,543 \\
\texttt{ai\_flagged\_post\_answers} & 94,217 \\
\texttt{post\_flag\_prevalence} & 5.50\% \\
\texttt{sensitivity\_sample\_exclude\_ai\_flagged} & 4,120,865 \\
\bottomrule
\end{tabular}
\end{table}

\begin{table}[t]
\centering
\caption{Sensitivity of H1--H4 coefficients to sample restrictions}
\label{tab:sensitivitysamples}
\scriptsize
\begin{tabular}{lccc}
\toprule
Restriction & $\beta_3$ (Post$\times$ICI) & $\beta_4$ (Post$\times$ECI) & $N$ \\
\midrule
Baseline (all retained) & $-0.153^{***}$ & $0.228^{***}$ & 4,215,082 \\
Exclude community-wiki & $-0.100^{**}$ & $0.224^{***}$ & 4,127,670 \\
Exclude self-answers & $-0.098^{**}$ & $0.222^{***}$ & 3,791,226 \\
Exclude AI-flagged & $-0.104^{**}$ & $0.231^{***}$ & 4,120,865 \\
Exclude all three & $-0.096^{*}$ & $0.218^{***}$ & 3,703,449 \\
\bottomrule
\multicolumn{4}{@{}p{\textwidth}@{}}{\scriptsize Note: All models include full
controls and fixed effects. ``Exclude AI-flagged'' removes answers selected
by a noisy text-classification proxy (precision 0.74; recall 0.81), not a
verified set of AI-authored answers. The row is therefore a perturbation
test, not evidence that the result holds in an AI-free sample.}
\end{tabular}
\end{table}

\section{AI-assisted text detector}
\label{app:detector}

Likely AI-assisted text was identified using a fine-tuned DeBERTa-v3-large classifier (checkpoint \texttt{microsoft/deberta-v3-large}, Hugging Face transformers 4.38.1) with a binary classification head. The model was trained on 18,000 labeled examples: 9,000 human-written answers sampled from January~2019--October~2022 and 9,000 synthetic answers generated by prompting GPT-3.5-Turbo (March~2023 snapshot) and GPT-4 (August~2023 snapshot) with 4,500 held-out Stack Overflow questions not present in the analytical sample. The positive class is defined as probability $\geq 0.72$, selected by maximizing Youden's $J$ on a 3,000-item validation partition. Training human texts are drawn exclusively from the pre-event window; classification targets are post-January~2023 answers.

Three trained annotators, blinded to detector scores, rated a stratified sample of 1,000 answers. The majority-vote reference is positive for 280 answers and negative for 720. Table~\ref{tab:detectorperf} reports performance. Cell counts are 227 true positives, 53 false negatives, 79 false positives, and 641 true negatives. Performance across sub-periods (2023--2025) showed no statistically significant decline in AUC (range 0.89--0.92; likelihood-ratio test $\chi^2(2)=3.1$, $p=0.21$). Of 1,712,539 post-period answers, 94,217 (5.50\%) exceeded the 0.72 threshold. This figure is an apparent classification rate, not an estimate of true AI-text prevalence. The flag defines sensitivity samples only and is not interpreted as verified authorship.

\begin{table}[t]
\centering
\caption{AI-assisted text detector performance ($n=1{,}000$ human-validation sample)}
\label{tab:detectorperf}
\scriptsize
\begin{tabular}{lc}
\toprule
Metric & Value (95\% CI) \\
\midrule
Sensitivity (recall) & 0.81 [0.76, 0.86] \\
Specificity & 0.89 [0.87, 0.91] \\
Precision & 0.74 [0.69, 0.79] \\
F1-score & 0.77 \\
AUC-ROC & 0.91 [0.89, 0.93] \\
Reference-positive rate & 28\% (280 / 1,000) \\
\bottomrule
\multicolumn{2}{@{}p{\textwidth}@{}}{\scriptsize Counts are 227, 53, 79,
and 641. Sensitivity is $227/280=0.81$, specificity is $641/720=0.89$,
and precision is $227/(227+79)=0.74$. The same precision is
$(0.81\times 0.28)/[(0.81\times 0.28)+(0.11\times 0.72)]=0.74$.
F1 is $2\times 227/(2\times 227+79+53)=0.77$. Wald intervals use those
cell counts.}
\end{tabular}
\end{table}

\section{Notation and coding anchors}
\label{app:notation}

Table~\ref{tab:notation} records principal indices and model notation. Table~\ref{tab:codinganchors} reports coding anchors for informational and epistemic contribution.

\begin{table}[t]
\centering
\caption{Principal indices and model notation}
\label{tab:notation}
\scriptsize
\begin{tabular}{@{}p{0.16\textwidth}p{0.46\textwidth}p{0.26\textwidth}@{}}
\toprule
Symbol & Standardized meaning & Primary use \\
\midrule
$a$ & Answer; the focal answer-level observation & H1--H4 \\
$q$; $q(a)$ & Question containing answer $a$ & ASI, complexity \\
$u$; $u(a)$ & Contributor; author of answer $a$ & H5/H7 panel \\
$t$ & Calendar month & Post indicator, FEs \\
$k$ & Event month relative to focal event & Event study \\
$g$ & Contribution function $\in\{\mathrm{I},\mathrm{E},\mathrm{C}\}$ & Theory only \\
$s(a)$ & Fixed-effect stratum: primary-tag $\times$ question-creation year & Eq.~\eqref{eq:did} \\
$Y_{at}$; $Y_{ut}$ & Answer-level; contributor-month outcome & Outcomes \\
$\mathrm{ICI}_a$; $\mathrm{ECI}_a$ & Informational; epistemic index & Valuation \\
$\mathrm{ASI}_q$ & Expert-rated AI answerability of question $q$ & Exposure proxy \\
$\mathrm{Post}_t$ & Post-event month indicator & Change specs \\
$\mathrm{Status}^{\mathrm{pre}}_u$ & Standardized log pre-event reputation & Adaptation \\
$\mathbf{X}_{at}$; $\mathbf{Z}_{ut}$ & Answer-level; contributor-month controls & Covariates \\
$\mathrm{Volatility}_{\operatorname{tag}(q),t}$ & Tag-level volatility & H6 \\
$\mathrm{RV}_{g,t}(\mathrm{ASI}_q)$ & Return to function $g$ & Theory \\
$\beta$; $\beta_k$; $\beta_1$--$\beta_4$; $\rho$ & Parameters & Estimands \\
$\theta_{\mathrm{I}}$; $\theta_{\mathrm{E}}$ & Triple-interaction coefficients & H3, H4 \\
$\boldsymbol{\gamma}$; $\boldsymbol{\kappa}$ & Control coefficient vectors & Controls \\
$\alpha_{s(a)}$; $\delta_t$ & Stratum; month fixed effects & Baseline \\
$\mu_{u(a)}$; $\lambda_{\operatorname{tag}(q(a))}$; $\eta_u$ & Contributor; tag; panel FEs & Valuation/adaptation \\
$\varepsilon_{at}$; $\varepsilon_{ut}$ & Disturbances & Errors \\
$\mathrm{EventTime}_t$; $\mathbb{1}\{\cdot\}$ & Event month; indicator & Event study \\
\bottomrule
\end{tabular}
\end{table}

\begin{table*}[t]
\centering
\caption{Coding anchors for informational and epistemic contribution}
\label{tab:codinganchors}
\scriptsize
\begin{tabular}{p{0.18\textwidth}p{0.34\textwidth}p{0.34\textwidth}}
\toprule
Pattern & Coding anchor & Illustrative answer form \\
\midrule
High ICI, low ECI & Supplies a direct command, code fragment, or standard
procedure without visible tests, assumptions, or limits & ``Use function
X with argument Y.'' \\
Low ICI, high ECI & Diagnoses uncertainty, compares explanations, or
specifies tests and boundary conditions without offering a complete
drop-in solution & ``This depends on the runtime version; test A before
choosing between B and C.'' \\
High ICI, high ECI & Supplies a usable solution and makes its warrant,
scope, and failure conditions inspectable & ``Use X; the example verifies
condition Y, but it fails under Z and requires version V.'' \\
Low ICI, low ECI & Provides neither a usable solution nor inspectable
reasoning & Unsupported assertion, social comment, or restatement of the
question \\
\bottomrule
\end{tabular}
\end{table*}

\section{Identification and robustness tables}
\label{app:identification}

Table~\ref{tab:secular_control} records that a $\mathrm{Post}_t\times\mathrm{PlatformAnswerVolume}_t$ term cannot be separately estimated alongside saturated month fixed effects. Table~\ref{tab:twoevent} reports the two-event piecewise specification within LLM arrival. Table~\ref{tab:prewindow} reports eight pre-window ASI interactions (joint $F(8)=1.14$, $p=0.33$). Table~\ref{tab:timeplacebo} reports placebo-in-time tests. Table~\ref{tab:lowcopilot} reports H1--H4 outside the first code-LLM languages ($N=1{,}475{,}279$). Table~\ref{tab:tagslope} adds a pre-period tag answer-volume slope control.

\begin{table}[t]
\centering
\caption{Platform-wide secular decline and identification}
\label{tab:secular_control}
\scriptsize
\setlength{\tabcolsep}{3pt}
\begin{tabular}{p{0.28\textwidth}p{0.20\textwidth}p{0.20\textwidth}p{0.24\textwidth}}
\toprule
Specification & Post$\times$ICI & Post$\times$ECI & Treatment of secular decline \\
\midrule
Baseline OLS & $-0.153$ & $0.228$ & Saturated calendar-month fixed effects \\
Add Post $\times$ platform-wide monthly answer volume & Not separately identified & Not separately identified & Collinear with calendar-month fixed effects \\
\bottomrule
\end{tabular}
\end{table}

\begin{table}[t]
\centering
\caption{Two-event piecewise specification within LLM arrival}
\label{tab:twoevent}
\scriptsize
\begin{tabular}{p{0.34\textwidth}p{0.28\textwidth}p{0.28\textwidth}}
\toprule
Window & Joint $F$-test ($p$) & Interpretation \\
\midrule
Code-LLM leads ($k^C=-6$ to $-1$) & $F(6)=0.71$ ($p=0.64$) & Joint test does not reject \\
Code-LLM lags ($k^C=0$ to $+5$) & $F(6)=1.27$ ($p=0.27$) & Short window; test does not reject \\
Leads before public LLM arrival ($k^G=-6$ to $-2$) & Same five-lead test as the event study & Not a second test \\
After public LLM arrival ($k^G=0$ to $+12$) & $F(13)=4.38$ ($p<0.001$) & Jointly significant association \\
\bottomrule
\end{tabular}
\end{table}

\begin{table}[t]
\centering
\caption{Full pre-window $\mathrm{ASI}$ interactions
(net vote score)}
\label{tab:prewindow}
\scriptsize
\begin{tabular}{lcc}
\toprule
Bin, relative to November 2022 & $\beta$ & SE \\
\midrule
January--June 2019 & 0.014 & 0.046 \\
July--December 2019 & $-0.021$ & 0.041 \\
January--June 2020 & 0.008 & 0.037 \\
July--December 2020 & $-0.033$ & 0.034 \\
January--June 2021 & 0.011 & 0.032 \\
July--December 2021 & $-0.027$ & 0.031 \\
January--May 2022 & $-0.019$ & 0.033 \\
June--October 2022 & $-0.046$ & 0.035 \\
\midrule
Joint test of the eight bins & \multicolumn{2}{c}{$F(8)=1.14$, $p=0.33$} \\
\bottomrule
\multicolumn{3}{@{}p{\textwidth}@{}}{\scriptsize The June--October 2022 bin
overlaps early code-LLM diffusion and is the least quiet pre-arrival bin.}
\end{tabular}
\end{table}

\begin{table}[t]
\centering
\caption{Placebo-in-time tests, pre-arrival sample only}
\label{tab:timeplacebo}
\scriptsize
\begin{tabular}{lcccc}
\toprule
Pseudo-event & Post$\times$ICI & SE & Post$\times$ECI & SE \\
\midrule
30 November 2019 & $-0.016$ & 0.047 & 0.024 & 0.061 \\
30 November 2020 & $-0.028$ & 0.043 & 0.019 & 0.056 \\
30 June 2021 & $-0.034$ & 0.041 & 0.038 & 0.053 \\
\midrule
LLM arrival (30 November 2022) & $-0.153$ & 0.036 & 0.228 & 0.041 \\
\bottomrule
\multicolumn{5}{@{}p{\textwidth}@{}}{\scriptsize The last row is the
reported arrival estimate. Each placebo coefficient is smaller than twice
its printed standard error. The largest absolute ratio of coefficient to
standard error is $0.83$. No stored $p$-value is reported for these
cells. The printed standard errors do not indicate a departure at $0.05$,
and the comparison does not prove that the interactions are unique to
late 2022.}
\end{tabular}
\end{table}

\begin{table}[t]
\centering
\caption{H1--H4 estimates outside the first code-LLM languages}
\label{tab:lowcopilot}
\scriptsize
\begin{tabular}{lccc}
\toprule
Coefficient & Reported, full sample & Outside first code-LLM languages & SE \\
\midrule
Post$\times$ICI & $-0.153$ & $-0.086$ & 0.049 \\
Post$\times$ECI & $0.228$ & $0.191$ & 0.064 \\
Post$\times$ASI$\times$ICI & $-0.139$ & $-0.117$ & 0.055 \\
Post$\times$ASI$\times$ECI & $0.176$ & $0.154$ & 0.058 \\
\bottomrule
\multicolumn{4}{@{}p{\textwidth}@{}}{\scriptsize Subsample
$N=1{,}475{,}279$.}
\end{tabular}
\end{table}

\begin{table}[t]
\centering
\caption{Control for pre-period tag answer-volume slopes}
\label{tab:tagslope}
\scriptsize
\begin{tabular}{lcc}
\toprule
Coefficient & Reported & With tag slope \\
\midrule
Post$\times$ICI & $-0.153$ (0.036) & $-0.087$ (0.039) \\
Post$\times$ECI & $0.228$ (0.041) & $0.194$ (0.053) \\
Post$\times$ASI$\times$ICI & $-0.139$ (0.045) & $-0.121$ (0.044) \\
Post$\times$ASI$\times$ECI & $0.176$ (0.052) & $0.149$ (0.046) \\
\bottomrule
\multicolumn{3}{@{}p{\textwidth}@{}}{\scriptsize Parentheses are standard
errors.}
\end{tabular}
\end{table}

\section{Inference, diagnostics, and multiplicity}
\label{app:inference}

Table~\ref{tab:pvalues} reports four-test and eight-test Bonferroni values for the primary net-vote family. Table~\ref{tab:family12} reports the twelve-test Bonferroni family. Table~\ref{tab:diagnostics} summarizes primary model diagnostics. Table~\ref{tab:vif} reports generalized VIFs.

\begin{table}[t]
\centering
\caption{Primary net-vote $p$-values with retrospective multiplicity corrections}
\label{tab:pvalues}
\scriptsize
\begin{tabular}{lccc}
\toprule
H & Raw $p$ & Bonferroni, four tests & Bonferroni, eight-test sensitivity \\
\midrule
H1 & 0.00002 & 0.00008 & 0.00016 \\
H2 & $<$0.0001 & $<$0.0004 & $<$0.0008 \\
H3 & 0.0020 & 0.0080 & 0.016 \\
H4 & 0.00071 & 0.0029 & 0.0057 \\
\bottomrule
\end{tabular}
\end{table}

\begin{table}[t]
\centering
\caption{Twelve-test Bonferroni values computed from printed estimates}
\label{tab:family12}
\scriptsize
\begin{tabular}{llcc}
\toprule
Coefficient & Source $p$ & $\times 12$ & Below 0.05 \\
\midrule
Post$\times$ICI, net votes & 0.00002 & 0.00024 & Yes \\
Post$\times$ECI, net votes & $<$0.0001 & $<$0.0012 & Yes \\
Post$\times$ASI$\times$ICI, net votes & 0.0020 & 0.024 & Yes \\
Post$\times$ASI$\times$ECI, net votes & 0.00071 & 0.0086 & Yes \\
Post$\times$ICI, acceptance & 0.031 & 0.375 & No \\
Post$\times$ECI, acceptance & 0.0003 & 0.004 & Yes \\
Post$\times$ASI$\times$ICI, acceptance & 0.006 & 0.075 & No \\
Post$\times$ASI$\times$ECI, acceptance & 0.0006 & 0.008 & Yes \\
H5, ECI share & $<$0.0001 & $<$0.0012 & Yes \\
H7, selected ASI & $<$0.0001 & $<$0.0012 & Yes \\
H6, volatility & 0.034 & 0.408 & No \\
Post$\times$status, count & 0.0014 & 0.017 & Yes \\
\bottomrule
\multicolumn{4}{@{}p{\textwidth}@{}}{\scriptsize Acceptance and count
$p$-values are two-sided Wald tests computed from the reported
coefficients and standard errors ($z=\hat\beta/\mathrm{SE}$).
H1, H3, and H4 net-vote $p$-values are two-sided normal probabilities
for the printed coefficient divided by its standard error.
H1 uses $t=-0.153/0.036=-4.25$, reported as $0.00002$; twelve times
that reported value is $0.00024$. H3 uses $t=-0.139/0.045=-3.09$
($p=0.0020$). H4 uses $t=0.176/0.052=3.38$ ($p=0.00071$).}
\end{tabular}
\end{table}

\begin{table*}[t]
\centering
\caption{Primary model diagnostics}
\label{tab:diagnostics}
\scriptsize
\setlength{\tabcolsep}{3pt}
\begin{tabular}{p{0.22\textwidth}p{0.17\textwidth}p{0.17\textwidth}p{0.17\textwidth}p{0.17\textwidth}}
\toprule
Diagnostic & Net vote score (OLS) & Acceptance (logit) & ECI share (frac.\ resp.) & Exit hazard (cloglog) \\
\midrule
Estimator & OLS & Logit & Fractional response & Complementary log--log \\
Clustering level & Tag (14,502) & Contributor (354,209) & Contributor & Contributor \\
Outcome mean / variance & 3.14 / 70.9 & 0.42 / 0.24 & 0.38 / 0.06 & Event rate 3.2\% \\
Distribution / zero share & Median 1; zero share 38.2\%; min $-12$, max 1,405 & 0 to 1 & 0 to 1 & 0 to 1 \\
Robust variance estimator & Tag-clustered HC & Contributor-clustered & Contributor-clustered & Contributor-clustered \\
Heteroskedasticity & White test $p<0.001$; tag-clustered robust SE & --- & --- & --- \\
Convergence & Closed-form OLS solution & No perfect prediction; Firth check stable & Converged & Converged \\
Link/specification test & RESET-type $p=0.26$ & Hosmer--Lemeshow $p=0.41$ & RESET-type $p=0.31$ & Time-varying coeff.\ $p=0.34$ \\
Influence / cluster sensitivity & Leave-one-tag-out maximum coefficient change $=0.019$ & --- & --- & --- \\
Calibration (Brier; slope) & --- & 0.21; 0.97 & --- & --- \\
Discrimination & --- & AUC = 0.71 & --- & C-index = 0.74 (attrition) \\
Omnibus Wald test & $p<0.001$ & $p<0.001$ & $p<0.001$ & $p<0.001$ \\
\bottomrule
\end{tabular}
\end{table*}

\begin{table}[t]
\centering
\caption{Generalized VIFs for the valuation specification and the separate supply specification (centered)}
\label{tab:vif}
\scriptsize
\begin{tabular}{lccc}
\toprule
Term & Df & GVIF & GVIF$^{1/(2\,\mathrm{Df})}$ \\
\midrule
\multicolumn{4}{l}{\textit{Valuation specification}} \\
$\mathrm{Post}_t$ & 1 & 1.02 & 1.01 \\
$\mathrm{ICI}_a$ & 1 & 1.18 & 1.09 \\
$\mathrm{ECI}_a$ & 1 & 1.24 & 1.11 \\
$\mathrm{ASI}_{q(a)}$ & 1 & 1.15 & 1.07 \\
$\mathrm{Post}_t\times\mathrm{ICI}_a$ & 1 & 2.31 & 1.52 \\
$\mathrm{Post}_t\times\mathrm{ECI}_a$ & 1 & 2.48 & 1.57 \\
$\mathrm{Post}_t\times\mathrm{ASI}_{q(a)}$ & 1 & 2.12 & 1.46 \\
$\mathrm{Post}_t\times\mathrm{ASI}_{q(a)}\times\mathrm{ICI}_a$ & 1 & 3.85 & 1.96 \\
$\mathrm{Post}_t\times\mathrm{ASI}_{q(a)}\times\mathrm{ECI}_a$ & 1 & 3.92 & 1.98 \\
\midrule
\multicolumn{4}{l}{\textit{Supply specification}} \\
ICI$\times$Post$\times$Supply & 1 & 3.84 & 1.96 \\
ECI$\times$Post$\times$Supply & 1 & 5.29 & 2.30 \\
\bottomrule
\end{tabular}
\par\smallskip
{\scriptsize Note: Df is the number of model-matrix columns assigned to
the term. Each displayed continuous interaction contributes one
column, so the adjusted value is the square root of the reported GVIF.
The two panels are separate specifications.}
\end{table}

\section{Measurement and ASI validation}
\label{app:measurement_ext}

Table~\ref{tab:modelasi} reports agreement between expert ASI and a frozen pre-30~November~2022 model checkpoint on the 600-question expert panel. Correctness rises from $0.19$ in the lowest ASI decile to $0.67$ in the highest; ASI correlates $0.58$ with correctness and $0.51$ with usability. Estimating H3 and H4 on the $N=1{,}380$ attached answers yields $\mathrm{Post}\times\mathrm{Correct}\times\mathrm{ICI}=-0.11$ (SE $0.09$) and $\mathrm{Post}\times\mathrm{Correct}\times\mathrm{ECI}=0.15$ (SE $0.10$), same sign as $-0.139$ and $0.176$ but too imprecise to confirm the answerability interpretation alone. Raising the AI-flag threshold to precision $0.85$ yields Post$\times$ICI $=-0.100$ (SE $0.037$) and Post$\times$ECI $=0.224$ (SE $0.050$), with $N=4{,}153{,}400$.

\begin{table}[t]
\centering
\caption{Agreement between expert ASI and a frozen model
($n=600$ questions)}
\label{tab:modelasi}
\scriptsize
\begin{tabular}{lc}
\toprule
Comparison & Estimate \\
\midrule
ASI with model correctness ($r$) & 0.58 \\
ASI with model usability ($r$) & 0.51 \\
ASI rank with correctness (Spearman) & 0.54 \\
Correctness, ASI decile 1 / decile 10 & 0.19 / 0.67 \\
Usability, ASI decile 1 / decile 10 & 0.24 / 0.61 \\
\bottomrule
\multicolumn{2}{@{}p{\textwidth}@{}}{\scriptsize The 600-question sample is
the expert panel used for ASI validation.}
\end{tabular}
\end{table}

\section{Additional results and sensitivity analyses}
\label{app:extra_results}

Table~\ref{tab:outcomerobust} reports top-1\% trim and ordered outcome specifications. Table~\ref{tab:attrition} and Table~\ref{tab:ipw} document attrition and IPW adaptation estimates. Table~\ref{tab:placebo} reports a placebo outcome. Table~\ref{tab:supply} reports an exploratory informational-supply association (not causal mediation). Table~\ref{tab:exitbounds} reports person-weighted exit scenarios. Table~\ref{tab:deleted} reports hypothetical deleted-answer bounds; these are not estimated regression coefficients and cannot recover deleted posts.

\begin{table*}[t]
\centering
\caption{Robustness of the primary interactions to alternative outcome specifications}
\label{tab:outcomerobust}
\scriptsize
\setlength{\tabcolsep}{4pt}
\begin{tabular}{llccc}
\toprule
Hypothesis & Variable & Baseline (OLS) & Trim top 1\% (OLS) & Ordered outcome (logit) \\
\midrule
H1 & Post$\times$ICI & $-0.153$ (0.036)$^{***}$ & $-0.141$ (0.038)$^{***}$ & $-0.412$ (0.098)$^{***}$ \\
H2 & Post$\times$ECI & $0.228$ (0.041)$^{***}$ & $0.215$ (0.043)$^{***}$ & $0.615$ (0.110)$^{***}$ \\
H3 & Post$\times$ICI$\times$ASI & $-0.139$ (0.045)$^{**}$ & $-0.125$ (0.048)$^{**}$ & $-0.380$ (0.125)$^{**}$ \\
H4 & Post$\times$ECI$\times$ASI & $0.176$ (0.052)$^{***}$ & $0.162$ (0.055)$^{**}$ & $0.475$ (0.140)$^{***}$ \\
\midrule
& Observations & 4,215,082 & 4,172,931 & 4,215,082 \\
& Adjusted $R^2$ / pseudo $R^2$ & 0.142 & 0.138 & 0.095 \\
\bottomrule
\multicolumn{5}{@{}p{\textwidth}@{}}{\scriptsize Standard errors in
parentheses are clustered at the tag level. $^{***}p<0.001$,
$^{**}p<0.01$, from the normal tail of the printed coefficient divided
by its standard error. Baseline cells repeat the net-vote column of
Table~\ref{tab:epistemic}. The trim drops the top 1\% of net-vote
answers. The ordered outcome uses a five-category vote binning. Ordered
coefficients are on the latent index.}
\end{tabular}
\end{table*}

\begin{table}[t]
\centering
\caption{Pre-event characteristics of retained versus exiting contributors}
\label{tab:attrition}
\scriptsize
\begin{tabular}{lcc}
\toprule
Pre-event characteristic & Retained ($n=218{,}447$) & Exited ($n=135{,}762$) \\
\midrule
Median pre-event reputation & 410 & 185 \\
Mean monthly answers (pre) & 2.8 & 1.1 \\
Mean ECI share & 0.38 & 0.31 \\
Mean selected ASI & 0.55 & 0.61 \\
Mean tag entropy & 1.9 & 1.2 \\
\bottomrule
\end{tabular}
\end{table}

\begin{table}[t]
\centering
\caption{Attrition-adjusted adaptation estimates}
\label{tab:ipw}
\scriptsize
\begin{tabular}{lccc}
\toprule
Outcome & Unweighted & IPW & Cohort (explicit zeros) \\
\midrule
ECI share & $0.058^{***}$ & $0.051^{***}$ & $0.047^{**}$ \\
& (0.012) & (0.014) & (0.015) \\
Selected ASI & $-0.082^{***}$ & $-0.074^{***}$ & $-0.069^{**}$ \\
& (0.019) & (0.022) & (0.024) \\
Complexity & $0.045^{**}$ & $0.039^{*}$ & $0.036^{*}$ \\
& (0.015) & (0.017) & (0.018) \\
\bottomrule
\multicolumn{4}{@{}p{\textwidth}@{}}{\scriptsize $^{*}p<0.05$, $^{**}p<0.01$, $^{***}p<0.001$.
The cohort epistemic-share estimate is $0.047/0.015=3.13$ (two-sided $p=0.0017$), so it takes two stars.
The cohort specification retains all 354,209 pre-period contributors.}
\end{tabular}
\end{table}

\begin{table}[t]
\centering
\caption{Placebo test: pre-period edit volume}
\label{tab:placebo}
\scriptsize
\begin{tabular}{lccc}
\toprule
Outcome & Interaction & Estimate & $p$-value \\
\midrule
Pre-period answer edit volume & Post$\times$ASI & 0.004 & 0.612 \\
\bottomrule
\multicolumn{4}{@{}p{\textwidth}@{}}{\scriptsize Note: The coefficient is
interpreted as a placebo differential association, not as proof of a
structural zero.}
\end{tabular}
\end{table}

\begin{table}[t]
\centering
\caption{Informational-supply association}
\label{tab:supply}
\scriptsize
\begin{tabular}{lc}
\toprule
Quantity & Estimate \\
\midrule
Post-period change in supply, high-ASI tags (log points) & 0.17 \\
Post-period change in supply, low-ASI tags (log points) & 0.04 \\
Supply$\times$ECI & 0.092 (SE 0.040) \\
Post$\times$ICI after supply control & $-0.077$ (SE 0.038) \\
Post$\times$ECI after supply control & 0.171 (SE 0.052) \\
Post$\times$ASI$\times$ECI after supply control & 0.128 (SE 0.046) \\
Informational interquartile contrast after supply control & $-0.026$ \\
Share of the $-0.153$ coefficient remaining & 0.503 \\
Share of the $0.228$ coefficient remaining & 0.75 \\
\bottomrule
\multicolumn{2}{@{}p{\textwidth}@{}}{\scriptsize High- and low-ASI tags are
the top and bottom quartiles of mean pre-period ASI. The supply term is
not interpreted as a causal mediator.}
\end{tabular}
\end{table}

\begin{table}[t]
\centering
\caption{Person-weighted scenarios for the status interactions}
\label{tab:exitbounds}
\scriptsize
\begin{tabular}{lcc}
\toprule
Assumption for the 135,762 exiters & H5, ECI share & H7, selected ASI \\
\midrule
Same shift as stayers (reported $0.058$, $-0.082$) & 0.058 & $-0.082$ \\
Half the stayer shift & 0.047 & $-0.066$ \\
No shift & 0.036 & $-0.051$ \\
Opposite shift of the same size & 0.014 & $-0.019$ \\
\bottomrule
\multicolumn{3}{@{}p{\textwidth}@{}}{\scriptsize Weights are $0.617$ on
stayers and $0.383$ on exiters. Calculated from the reported
coefficients and the attrition counts. Not an estimate of exiter
behavior.}
\end{tabular}
\end{table}

\begin{table*}[t]
\centering
\caption{Hypothetical sensitivity bounds under assumed deleted-answer scenarios}
\label{tab:deleted}
\scriptsize
\setlength{\tabcolsep}{3pt}
\begin{tabular}{lcccccc}
\toprule
Scenario & \makecell{Post-period\\deletion} & \makecell{Pre-period\\deletion} & \makecell{Vote penalty\\high ICI} & \makecell{Vote penalty\\high ECI} & \makecell{Hypothetical\\Post$\times$ICI bound} & \makecell{Hypothetical\\Post$\times$ECI bound} \\
\midrule
1 Baseline & 8\% & 3\% & $-1.8$ & $-0.4$ & $-0.119$ & $0.211$ \\
2 Conservative & 5\% & 2\% & $-1.8$ & $-0.4$ & $-0.131$ & $0.218$ \\
3 Aggressive & 10\% & 5\% & $-1.8$ & $-0.4$ & $-0.105$ & $0.202$ \\
4 High penalty & 8\% & 3\% & $-2.5$ & $-0.6$ & $-0.108$ & $0.195$ \\
5 Worst case & 10\% & 5\% & $-2.5$ & $-0.6$ & $-0.092$ & $0.184$ \\
6 Minimal impact & 2\% & 1\% & $-1.0$ & $-0.2$ & $-0.142$ & $0.225$ \\
\bottomrule
\multicolumn{7}{@{}p{\textwidth}@{}}{\scriptsize These values are
hypothetical bounds from applying assumed vote penalties to assumed
deletion rates. They are not estimated regression coefficients, and they
cannot be reproduced from the exported covariance matrices. The signs
match the reported $-0.153$ and $0.228$ in every row. The bounds
illustrate directional robustness under those assumptions. They do not
recover the deleted posts.}
\end{tabular}
\end{table*}
